\article{函数域上 $\GL_2$ 的自守形式（依照 V. G. Drinfeld）}{Automorphic Forms on $\GL_2$ over Function Fields (after V. G. Drinfeld)}{G. Harder 与 D. A. Kazhdan}{2}{357--379}
\paragraph*{引言（Introduction）}
本文报告 V. G. Drinfeld 关于函数域（function fields）上自守形式（automorphic forms）的尚未发表的工作。本文将描述 shtuka 概形（scheme of shtuka），说明如何用 Selberg 迹公式（Selberg trace formula）计算 Frobenius 在其上同调（cohomology）上的迹，并由此证明 Ramanujan 猜想（Ramanujan conjecture）。我们不讨论与互反律（reciprocity law）有关的、精细得多的问题；这个问题涉及上述模概形（moduli scheme）的紧化（compactification）。

本文给出证明概要，并附以参考文献和必要提示，供读者补全细节。
\footnote{AMS（MOS）主题分类（1970）：主要分类 10D15、14F05、14H10。原文版权：\textcopyright\ 1979，美国数学会。}
\origsection{1}{$FH$ 层的模空间（The moduli space of $FH$-sheaves）}
\paragraph*{1.0.}
设 $X/\F_q$ 为光滑射影曲线（smooth projective curve），记 $\bar X=X\times_{\F_q}\bar\F_q$。Drinfeld 在 $\bar X$ 上引入了称为 shtuka 的对象。它们是 $\bar X$ 上带有附加结构的向量丛（vector bundles），其中的附加结构由 Frobenius 与某种 Hecke 修正（Hecke modifications）给出。因此，我们建议称它们为 $FH$ 层（$FH$-sheaves）。
\origsubsection{1.1}{向量丛（Vector bundles）}
设 $X/\F_q$ 为给定曲线，$K/\F_q$ 为其亚纯函数域（field of meromorphic functions）。记
\[
\bar X=X\times_{\F_q}\bar\F_q,\qquad K_\infty=K\cdot\bar\F_q.
\]
$K$（分别为 $K_\infty$）的阿代尔空间（space of adeles）记为 $A$（分别为 $A_\infty$）。$X/\F_q$ 的几何点（geometric points）记为 $v,w\in X(\bar\F_q)$。对任意 $v\in X(\bar\F_q)$，$K_\infty$ 关于 $v$ 的完备化（completion）记为 $K_v$，其整数环（ring of integers）记为 $\mathcal O_v$。

域 $K/\F_q$ 的赋值（valuations）记为 $\mathfrak p,\ldots$；$K$ 在 $\mathfrak p$ 处的完备化记为 $K_{\mathfrak p}$，整数环为 $\mathcal O_{\mathfrak p}$。若 $v\in X(\bar\F_q)$ 在 $K$ 上诱导 $\mathfrak p$，则写作 $v\to\mathfrak p$。位于同一 $\mathfrak p$ 之上的点 $v$ 构成 Frobenius 在 $X(\bar\F_q)$ 上作用的一个轨道（orbit）。

$\bar X$ 上秩为 $d$ 的向量丛，是结构层（structure sheaf）$\mathcal O_{\bar X}$ 上秩为 $d$ 的局部自由层（locally free sheaf）。向量丛也可用各处的格来描述。取带基 $e_1,\ldots,e_d$ 的向量空间 $V/K_\infty$。设有一族 $\mathcal O_v$ 格（lattices）
\[
M=\{M_v\}_{v\in X(\bar\F_q)},\qquad M_v\subset V\otimes K_v,
\qquad M_v=\bigoplus_{i=1}^d\mathcal O_ve_i\quad\text{对几乎所有 }v.
\]
% 原书第358页
这族格确定一个局部自由层 $E/\bar X$：对于开子集 $\bar U\subset\bar X$，令
\[
\Gamma(\bar U,E)=\{x\in V\mid x\in M_v\text{，对所有 }v\in\bar U\}.
\]
反过来，每个向量丛都可由这样的格族得到。

取平凡格族（trivial family）$M_0$，即对所有 $v\in X(\bar\F_q)$ 令
$M_{0,v}=\mathcal O_ve_1\oplus\cdots\oplus\mathcal O_ve_d$。阿代尔 $x\in\GL_d(A_\infty)$ 给出新的格族 $xM_0=\{x_vM_{0,v}\}$。容易验证，下式给出一个双射（bijection）：
\[
\GL_d(K_\infty)\backslash\GL_d(A_\infty)/\mathcal K_\infty
\xrightarrow{\sim}\{\bar X\text{ 上秩为 }d\text{ 的向量丛的同构类}\},
\qquad\mathcal K_\infty=\prod_v\GL_d(\mathcal O_v).
\]
\origsubsection{1.2}{向量丛的水平结构（Level structures on vector bundles）}
设 $D$ 为 $X/\F_q$ 上的正除子（positive divisor），即在 $\F_q$ 上有理的除子。按通常定义，$\mathcal O_{\bar X}(-D)$ 是满足 $\operatorname{div}(f)\geq D$ 的正则函数芽（germs of regular functions）所成的层，也就是定义 $D$ 的理想层（sheaf of ideals）。对任意向量丛，定义
$E(-D)=E\otimes\mathcal O_{\bar X}(-D)$。$E$ 上沿 $D$ 的水平结构（level structure）是一个 $\mathcal O_{\bar X}$ 同构
\[
\psi:E/E(-D)\xrightarrow{\sim}
\bigl(\mathcal O_{\bar X}/\mathcal O_{\bar X}(-D)\bigr)^d.
\]
沿 $D$ 带有水平结构的向量丛，其同构类由陪集空间（coset space）

\[
 \GL_d(K_\infty)\backslash\GL_d(A_\infty)/\mathcal K_\infty(D)
\]
的元素给出；这里 $\mathcal K_\infty(D)$ 是 $\mathcal K_\infty$ 中由 $x\equiv1\pmod D$ 定义的开同余子群（open congruence subgroup）。

我们特别关心 $d=1$ 与 $d=2$ 两种情形。带有沿 $D$ 的水平结构的秩二向量丛 $E/\bar X$ 称为稳定的（stable），如果对任意线子丛（line subbundle）$L\subset E$，都有
\[
\deg L<\frac12(\deg E+\deg D).
\]
Mumford 理论表明，沿 $D$ 带有水平结构的稳定丛所构成的模函子（functor）具有精细模概形（fine moduli scheme）$\mathcal M_0\to\Spec(\F_q)$（参见 [6]）。若对于 $\nu\in\Z$，以 $\mathcal M_D^{(\nu)}$ 表示满足 $\deg\det E=\nu$ 的稳定向量丛，则 $\mathcal M_D^{(\nu)}$ 是 $\F_q$ 上的光滑拟射影概形（smooth quasi-projective schemes）。
\origsubsection{1.3}{$d=2$ 时向量丛的修正（Modifications of vector bundles）}
设 $E/\bar X$ 是向量丛，$v\in X(\bar\F_q)$。我们要定义 $E$ 在 $v$ 处的修正。向量丛在 $v$ 处的纤维（fiber）
\[
E_v=E\otimes(\mathcal O_v/\mathfrak m_v)=E\otimes\bar\F_q(v)
\]
是 $\bar\F_q$ 上的向量空间。点 $\alpha_v\in\mathbf P^1(E_v)(\bar\F_q)$ 定义一个非零线性形式 $\alpha_v^*:E_v\to\bar\F_q$，它确定到相差一个非零标量。定义 $E(\alpha_v)$ 为 $E$ 中满足 $\alpha_v^*(s_v)=\alpha_v(s_v)=0$ 的截面芽（germs of sections）所成的层。这样得到 $\bar X$ 上的一个新向量丛，并有正合列（exact sequence）
\[
0\longrightarrow E(\alpha_v)\longrightarrow E\longrightarrow\bar\F_q(v)\longrightarrow0.
\]
也就是说，$E/E(\alpha_v)$ 是支撑集中于 $v$ 的层，它在 $v$ 处的纤维是剩余域（residue field）$\bar\F_q=\bar\F_q(v)$ 上的一维向量空间。我们称 $E(\alpha_v)$ 为 $E$ 在 $v$ 处（沿方向 $\alpha_v$）的下修正（lower modification）。

下面用格来解释。设有格族 $M=xM_0$，$x\in\GL_2(A_\infty)$。选取
\[
\alpha_v\in\GL_2(\mathcal O_v)
\begin{pmatrix}\pi_v&0\\0&1\end{pmatrix}\GL_2(\mathcal O_v),
\]
其中 $\pi_v$ 是一致化参数（uniformizing parameter）。置
$\boldsymbol\alpha_v=(\ldots,1,\ldots,\alpha_v,\ldots,1,\ldots)$，则由格族
$M(\alpha_v)=x\boldsymbol\alpha_v\cdot M_0$ 定义的丛显然给出 $M$ 的一个修正。
% 原书第359页
注意，集合
\[
\GL_2(\mathcal O_v)\begin{pmatrix}\pi_v&0\\0&1\end{pmatrix}\GL_2(\mathcal O_v)
\]
对右侧 $\GL_2(\mathcal O_v)$ 的作用取商后，所得集合恰可识别为 $\mathbf P^1(E_v)(\bar\F_q)$。这便给出了修正概念的格族解释。

若 $w\in X(\bar\F_q)$ 是另一点，则定义 $w$ 处的上修正（upper modifications）如下：对 $\alpha_w\in\mathbf P^1(E_w)(\bar\F_q)$，置
$E(\widetilde\alpha_w)=E(\alpha_w)\otimes\mathcal O_{\bar X}(w)$。此时有
$E\hookrightarrow E(\widetilde\alpha_w)\to\bar\F_q(w)$。用格族描述，就是
$M=xM_0\to x\widetilde{\boldsymbol\alpha}_wM_0$，其中
\[
\widetilde{\boldsymbol\alpha}_w=(\ldots,1,\ldots,\widetilde\alpha_w,\ldots,1,\ldots),\qquad
\widetilde\alpha_w\in\GL_2(\mathcal O_w)
\begin{pmatrix}\pi_w^{-1}&0\\0&1\end{pmatrix}\GL_2(\mathcal O_w).
\]
若 $v\ne w$ 为两个几何点，$\alpha_v\in\mathbf P^1(E_v)(\bar\F_q)$，$\alpha_w\in\mathbf P^1(E_w)(\bar\F_q)$，定义
$E(\alpha_v,\widetilde\alpha_w)=(E(\alpha_v))(\widetilde\alpha_w)$。注意
$\deg E(\alpha_v,\widetilde\alpha_w)=\deg E$。

此外，还有如下典范识别（canonical identification）
\[
E(\alpha_v,\widetilde\alpha_w)|_{\bar X-\{v\}-\{w\}}
=E|_{\bar X-\{v\}-\{w\}}.
\]
由 $E(\alpha_v,\widetilde\alpha_w)$ 的局部自由层定义，这一点显然成立。因此，若 $D$ 是 $X/\F_q$ 上的正除子，$v,w\notin\operatorname{supp}(D)$，则 $E$ 上沿 $D$ 的水平结构 $\psi$ 在 $E(\alpha_v,\widetilde\alpha_w)$ 上诱导一个水平结构，仍记为 $\psi$。

给出这样的双重修正（double modification）
$E\to E(\alpha_v,\widetilde\alpha_w)=E'$，等价于给出图式
\[
\begin{tikzcd}[column sep=large]
E\arrow[dr,"\phi_v"] & \\
&\mathcal E\\
E'\arrow[ur,"\phi_w"'] &
\end{tikzcd}
\qquad\mathcal E\text{ 是秩二局部自由层},
\]
其中 $\deg\mathcal E=\deg E+1=\deg E'+1$，且
$\mathcal E/\phi_v(E)$ 支撑集中于 $w$，$\mathcal E/\phi_w(E')$ 支撑集中于 $v$。只须取 $\mathcal E=E(\widetilde\alpha_w)$，并令 $\phi_v,\phi_w$ 为包含映射（inclusions）即可。
\origsubsection{1.4}{$FH$ 层及其模空间（The $FH$-sheaves and their moduli space）}
$\bar\F_q$ 的 Frobenius 自同构（Frobenius automorphism）$x\mapsto x^q$ 诱导映射 $\operatorname{Fr}:\Spec(\bar\F_q)\to\Spec(\bar\F_q)$，从而诱导
$\operatorname{Id}\times\operatorname{Fr}:\bar X=X\times_{\F_q}\Spec(\bar\F_q)\to\bar X$。

若 $E/\bar X$ 是向量丛，以 $E^\sigma$ 表示在上述映射下的拉回（pull-back），即
$E^\sigma=(\operatorname{Id}\times\operatorname{Fr})^*E$。若 $E$ 有沿 $D$ 的水平结构 $\psi$，便得到水平结构
\[
\psi^\sigma:E^\sigma/E^\sigma(-D)\longrightarrow
\bigl((\mathcal O_{\bar X}/\mathcal O_{\bar X}(-D))^2\bigr)^\sigma
=(\mathcal O_{\bar X}/\mathcal O_{\bar X}(-D))^2,
\]
因为 $(\mathcal O_{\bar X}/\mathcal O_{\bar X}(-D))^2
=(\mathcal O_X/\mathcal O_X(-D))^2\otimes_{\F_q}\bar\F_q$。
一组数据 $(E,\psi,\alpha_v,\alpha_w,\Phi)$ 称为一个 $FH$ 层，或 shtuka，如果满足：
\begin{enumerate}[label=(\arabic*)]
\item $D$ 是正有理除子（positive rational divisor），$v,w\in X(\bar\F_q)$，$v\ne w$ 且 $v,w\notin\operatorname{supp}(D)$；
\item $\alpha_v\in\mathbf P^1(E_v)(\bar\F_q)$，$\alpha_w\in\mathbf P^1(E_w)(\bar\F_q)$；
\item $\psi$ 是沿 $D$ 的水平结构；
\item $\Phi:(E^\sigma,\psi^\sigma)\xrightarrow{\sim}(E(\alpha_v,\widetilde\alpha_w),\psi)$。
\end{enumerate}
若固定 $v,w$ 两点，就称其为 $(v,w)$ 上的 $FH$ 层。
% 原书第360页
两个这样的 $FH$ 层称为同构（isomorphic），如果存在向量丛的同构，且它与全部附加结构相容。
\paragraph*{注（Remark）}
若 $(E,\psi,\alpha_v,\alpha_w,\Phi)$ 是 $FH$ 层，且 $(E,\psi)$ 足够稳定，例如对每个线子丛 $L\subset E$ 都有
\[
\deg L\leq\frac12(\deg E+\deg D)-4,
\]
则附加数据 $\alpha_v,\widetilde\alpha_w,v,w,\Phi$ 由 $(E,\psi)$ 唯一确定。这一点可直接由稳定性条件验证。

设已给定 $FH$ 层 $(E,\alpha_v,\widetilde\alpha_w,\Phi)$ 和 $X/\F_q$ 上的正有理除子 $D$。下面确定这个 $FH$ 层上所有可能的沿 $D$ 的水平结构。同构 $\Phi$ 诱导同构
$\Phi_D:E^\sigma/E^\sigma(-D)\to E/E(-D)$，水平结构映射
$\psi:E/E(-D)\to(\mathcal O_{\bar X}/\mathcal O_{\bar X}(-D))^2$ 必须满足
$\psi^{-1}\psi^\sigma=\Phi_D$。由 Lang 定理（Lang's theorem），这个方程总有解，且解集在
$\GL_2(\mathcal O_X/\mathcal O_X(-D))$ 的作用下构成主齐性空间（principal homogeneous space）。也就是说，所有其他解都是 $g\circ\psi$，其中
\[
g\in\GL_2(\mathcal O_X/\mathcal O_X(-D))=G(D).
\]

下面说明如何为 $FH$ 层的模函子构造粗模空间（coarse moduli space）。固定两个互异点 $v,w\in X(\bar\F_q)$，以及满足 $v,w\notin\operatorname{supp}(D)$ 的 $D$。对每个 $t\in\mathbf N$，令 $\mathcal M_{D,t}$ 为带沿 $D$ 的水平结构的向量丛 $E$ 的模空间，并要求每个线子丛 $L\subset E$ 都满足
\[
\deg L\leq\tfrac12(\deg E+\deg D)-t.
\]
这是上述 $\mathcal M_D$ 的开子概形（open subscheme）。这样的丛经修正得到的 $E(\alpha_v,\widetilde\alpha_w)$ 对应于 $\mathcal M_{D,t-1}$ 中的点。不难证明存在 $\bar\F_q$ 上的闭子概形（closed subscheme）
\[
\mathcal H_{v,w}\subset\mathcal M_{D,t}\times\mathcal M_{D,t-1},
\]
其几何点是成对的丛 $(E,E(\alpha_v,\widetilde\alpha_w))$。若 $\mathfrak p,\mathfrak q$ 是 $v,w$ 在 $K$ 上诱导的赋值，而 $\F_{q^d}$ 是 $\mathfrak p,\mathfrak q$ 的剩余域的合成域，则显然 $\mathcal H_{v,w}$ 定义在 $\F_{q^d}$ 上。

由前面的注可知，例如当 $t\geq4$ 时，修正的方向 $\alpha_v,\alpha_w$ 由 $(x,y)$ 唯一确定。此外，简单的无穷小论证（infinitesimal arguments）表明，$\mathcal H_{v,w}$ 在 $(x,y)$ 处的切空间（tangent space）满射到第一分量的切空间，而这个投影的核的维数为 $2$。

令 $\Gamma_{\operatorname{Fr}}$ 为几何 Frobenius 的图像（graph），即
$(x,x^{(q)})\subset\mathcal M_{D,t}\times\mathcal M_{D,t-1}$。那么显然
\[
M_{D,t}^{(v,w)}=\mathcal H_{v,w}\cap\Gamma_{\operatorname{Fr}}
\]
是 $\F_{q^d}$ 上的光滑二维概形，其几何点的第一分量恰好是带有唯一确定的 $FH$ 层结构的丛。

有限群 $G(D)$ 作用于 $M_{D,t}^{(v,w)}$。若 $D'\geq D$，则
$M_{D',t}^{(v,w)}$ 的某个开部分 $\widetilde M_{D',t}^{(v,w)}$ 有到 $M_{D,t}^{(v,w)}$ 的自然映射（当 $D$ 增大时，稳定性条件随之改变），且显然有
\[
\begin{tikzcd}[column sep=large]
\widetilde M_{D',t}^{(v,w)}\arrow[r,hook]\arrow[d]
&M_{D',t}^{(v,w)}\\
\widetilde M_{D',t}^{(v,w)}/\ker(G(D')\to G(D))\arrow[r,"\sim"]
&M_{D,t}^{(v,w)}.
\end{tikzcd}
\]
因此有 $M_{D,t'}^{(v,w)}\hookrightarrow M_{D',t}^{(v,w)}/\ker(G(D')\to G(D))$，并且对于某个 $t\geq4$，可定义
\[
M_D^{(v,w)}=\mathop{\operatorname{inj\,lim}}_{D'\geq D;\ v,w\notin\operatorname{supp}(D')}
M_{D',t}^{(v,w)}/\ker(G(D')\to G(D)).
\]
% 原书第361页
以上给出了下述定理的证明概要。
\begin{originaltheorem}{定理 1.4.1（Theorem 1.4.1）}
\textup{(i)} 带有沿 $D$ 的水平结构、位于 $(v,w)$ 上的 $FH$ 层函子，可由概形的归纳极限（inductive limit）$\{M_D^{(v,w)}\}/\F_{q^d}$ 表示。若固定行列式（determinant）的次数，则它是拟射影概形的极限，而这些拟射影概形是光滑概形对有限群作用所取的商。

\textup{(ii)} 令
$\Lambda_D=\bar X\times\bar X-\Delta-D\times\bar X-\bar X\times D$，则可构造归纳系统（inductive system）$M_D\to\Lambda_D$，其在 $(v,w)$ 上的纤维为 $M_D^{(v,w)}$。
\end{originaltheorem}
断言 (ii) 直接来自这样一个事实：也可以把 $v,w$ 作为可变参数来构造 $\mathcal H_{v,w}$。
\paragraph*{注（Remark）}
严格地说，这里使用的是相应的归纳系统；记作极限只是为了简化表述。后面还将讨论射影系统（projective system）$\operatorname{proj\,lim}_D M_D^{(v,w)}$。
\origsubsection{1.5}{$M_D^{(v,w)}$ 与 $M_D$ 的几何性质（Geometric properties）}
\paragraph*{1.5.1.}
先考虑 $d=1$。线丛（line bundle）$L/\bar X$ 有 $FH$ 层结构，当且仅当 $L^\sigma\to L(w-v)$，等价地，$L^{-1}\otimes L^\sigma=\mathcal O_{\bar X}(w-v)$。这说明，一维 $FH$ 层在 $\operatorname{Pic}_X(\F_q)$ 的作用下构成主齐性空间；这个群就是 $X$ 上线丛同构类组成的群。因此，$(v,w)$ 上一维 $FH$ 层的模空间是无限离散集（infinite discrete set）。

引入水平结构后，它成为 $X$ 上带有沿 $D$ 的有理水平结构的线丛群 $\operatorname{Pic}_{X,D}(\F_q)$ 的主齐性空间。

若固定 $D$，而让 $v,w$ 变化，便得到无限覆盖（infinite covering）$\widetilde P_D\to\Lambda_D$，其 Galois 群，也就是覆盖变换群（deck transformation group，原文德语 Decktransformationen），恰好是 $\operatorname{Pic}_{X,D}(\F_q)$。
\paragraph*{1.5.2.}
若 $d=2$，$(E,\psi)$ 是带水平结构的 $FH$ 层，则 $(\det E,\det\psi)$ 是带有水平结构的一维 $FH$ 层。由此得到映射
$\det:M_D\to\widetilde P_D$。
\origsubsection{1.6}{$FH$ 层的扩张（Extensions of $FH$-sheaves）}
现在假设 $v,w$ 不仅互异，而且不在 Frobenius 映射的同一轨道内。这等价于 $v,w$ 在 $K$ 上诱导不同赋值。

设有 $FH$ 层 $(E,\psi,\alpha_v,\alpha_w,\Phi)$，以及线子丛 $F\subset E$，使得
$\Phi:F^\sigma\to F'$，其中 $F'$ 是 $E(\alpha_v,\widetilde\alpha_w)$ 中与 $F$ 对应的线子丛（线子丛由其泛纤维（generic fiber）确定）。此时有且仅有以下两种可能。

(i) $\alpha_v$ 在纤维 $F\otimes\bar\F_q(v)=F_v\subset E_v$ 上非零。此时 $\alpha_w$ 在 $F_w$ 上为零，且 $F'=F(w-v)$，即 $F$ 本身是一维 $FH$ 层。在这种情形，商线丛（quotient line bundle）$E/F$ 定义在 $\F_q$ 上；这是 1.5.2 的推论。

(ii) $\alpha_v$ 在 $F_v$ 上为零，因而在 $F_w$ 上非零。此时 $F'=F=F^\sigma$，这个线丛定义在 $\F_q$ 上，而 $E/F$ 是一维 $FH$ 层。

在第一种情形，我们称这个 $FH$ 层为左 $FH$ 扩张（left $FH$-extension），把 $F$ 记为 $L$，写成
$0\to L\to E\to H\to0$。第二种情形称为右扩张（right extension），把 $F$ 记为 $H_1$，写成
$0\to H_1\to E\to L_1\to0$。下面研究一个 $FH$ 层何时能表示为扩张，以及这种表示有多少种。

考虑 $E$ 的泛纤维 $V/K_\infty$。$\Phi$ 诱导线性映射 $\Phi_K:V^\sigma\to V$。由于 $V^\sigma$ 与 $V$ 的底层加法群相同，这个映射诱导 $\Phi_K:V\to V$，但它只是 $\sigma^{-1}$ 线性的（$\sigma^{-1}$-linear），即
\[
\Phi_K(\lambda^\sigma\cdot x)=\lambda\Phi_K(x).
\]
% 原书第362页
因此，为使 $(E,\psi,\alpha_v,\alpha_w,\Phi)$ 是扩张，映射 $\Phi_K$ 必须有一维不变子空间（invariant subspace）$W\subset V$。于是可以选基向量 $e_1\in W$ 和 $\bar e_2\in V/W$，使得
$\Phi_K(e_1)=\lambda_1e_1$，$\Phi_K(\bar e_2)=\lambda_2\bar e_2$。
\begin{originaltheorem}{命题 1.6.1（Proposition 1.6.1）}
设 $\mathfrak p$（分别为 $\mathfrak q$）是 $v$（分别为 $w$）在 $K$ 上诱导的赋值，且 $\mathfrak p\ne\mathfrak q$。设
$\lambda_1,\lambda_2\in K\F_{q^n}=K_n$，并令
\footnote{译注：原书这一行将两个 $\lambda$ 的下标都印作 $1$；这里依前后文将第二个写为 $2$。}
\[
\mu_i=\operatorname{Norm}_{K_n/K}(\lambda_i)\in K,\qquad i=1,2.
\]
那么，在必要时交换 $\mu_1,\mu_2$ 的顺序后，有：
\begin{enumerate}[label=(\arabic*)]
\item $\mu_1$ 在 $\mathfrak p$ 处有零点，在 $\mathfrak q$ 处为单位；
\item $\mu_2$ 在 $\mathfrak q$ 处有极点，在 $\mathfrak p$ 处为单位；
\item 向量空间 $V$ 唯一分解为两个子空间之和 $V=W\oplus W'$，即 $W$ 有唯一的 $\Phi_K$ 不变补空间；
\item 若 $L$（分别为 $H_1$）是 $W$（分别为 $W'$）在 $E$ 中诱导的子丛，则 $(E,\psi,\alpha_v,\alpha_w,\Phi)$ 恰有一种写成左扩张的方式，也恰有一种写成右扩张的方式：
\[
0\to L\to E\to H\to0,\qquad
0\to H_1\to E\to L_1\to0.
\]
\end{enumerate}
\end{originaltheorem}
此命题可通过分析 $v,w$ 两处的局部情形证明（local analysis）。

本节最后讨论扩张的分类问题。给定 $L,H$ 或 $H_1,L_1$，如何描述 $FH$ 扩张
$0\to L\to E\to H\to0$ 与 $0\to H_1\to E_1\to L_1\to0$？

只讨论第一种情形。固定 $v,w$ 与除子 $D$，仍假设 $\mathfrak p\ne\mathfrak q$、$v,w\notin\operatorname{supp}D$。另外固定水平结构
\[
\psi_L:L/L(-D)\xrightarrow{\sim}\mathcal O_{\bar X}/\mathcal O_{\bar X}(-D),
\quad
\psi_H:H/H(-D)\xrightarrow{\sim}\mathcal O_{\bar X}/\mathcal O_{\bar X}(-D).
\]
\footnote{译注：原书第二个商写为 $H/L(-D)$；这里按水平结构的定义写为 $H/H(-D)$。}
以 $L$ 为子对象、$H$ 为商的扩张由
$H^1(\bar X,\Hom(H,L))$ 中的元素分类，而集合
\[
H^1(\bar X,\Hom(H,L)(-D))
=H^1(\bar X,H^{-1}\otimes L(-D))
\]
分类的是沿除子 $D$ 带有分裂（splitting）的扩张。

因此，对应于元素 $\xi\in H^1(\bar X,H^{-1}\otimes L(-D))$ 的向量丛 $E$，自动带有沿 $D$ 的水平结构。假设 $D$ 的次数相对于 $\deg(H^{-1}\otimes L)$ 很大。我们要找出哪些 $\xi$ 使所得向量丛带有一个（因而唯一的）$FH$ 结构。这等价于构造如下交换图\footnote{\textbf{译者注（本次校订）：}下图中间一行的中项应另记为 $\mathcal E=E(\widetilde\alpha_w)$。它的次数为 $\deg E+1$，不能仍与第一行的 $E$ 相同；这一点也可由该行的子丛 $L(w)$ 和商丛 $H$ 的次数相加看出。}
\[
\begin{tikzcd}[row sep=large]
0\arrow[r]&L\arrow[r]\arrow[dr,"i"]&E\arrow[r]\arrow[dr]&H\arrow[r]&0\\
0\arrow[r]&L(w)\arrow[r]&E\arrow[r]&H\arrow[r]&0\\
0\arrow[r]&L^\sigma\arrow[r]\arrow[ur,"\widetilde\Phi_L"]&E^\sigma\arrow[r]\arrow[ur]&H\arrow[r]&0.
\end{tikzcd}
\]

这里 $\widetilde\Phi_L$ 是复合映射
$L^\sigma\xrightarrow{\Phi_L}L(w-v)\hookrightarrow L(w)$。从而得到下图：
% 原书第363页
\[
\begin{tikzcd}[column sep=large,row sep=large]
H^1(\bar X,H^{-1}\otimes L(-D))\arrow[dr,"i^*"]&\\
&H^1(\bar X,H^{-1}\otimes L(-D+w))\\
H^1(\bar X,H^{-1}\otimes L^\sigma(-D))\arrow[ur,"\widetilde\Phi_L^*"']\arrow[d,equal]&\\
H^1(\bar X,H^{-1}\otimes L(-D))&
\end{tikzcd}
\]
要得到 $FH$ 层，必须寻找满足
\[
i^*(\xi)-\widetilde\Phi_L^*(\xi^\sigma)=0
\]
的类 $\xi$。
\begin{originaltheorem}{命题 1.6.2（Proposition 1.6.2）}
\textup{(1)} 这个方程定义一个 $\bar\F_q$ 上的光滑一维群概形（group scheme），它同构于 $\Ga$。

\textup{(2)} 它可定义在 $(L,\psi_L,\Phi_L)$ 与 $(H,\psi_H)$ 的定义域（field of definition）上。

\textup{(3)} 对每一对 $(L,\psi_L)$、$(H,\psi_H)$，得到两条仿射直线（affine lines）$G_D^{L,H}$、$G_D^{H,L}$，分别参数化带有沿 $D$ 的水平结构的右、左 $FH$ 扩张。\footnote{\textbf{译者注（本次校订）：}按1.6的定义，这里的“右、左”应依次为“左、右”：$G_D^{L,H}$ 参数化 $0\to L\to E\to H\to0$，$G_D^{H,L}$ 参数化以 $H$ 为子丛的扩张。}
\end{originaltheorem}
证明的主要步骤是说明上述方程的解空间是连通概形（connected scheme）。
\paragraph*{注（Remarks）}
1. 上述扩张构造所得的 $FH$ 层，带有一种特殊的水平结构，即与扩张适配的（adapted）水平结构：
$\psi:E/E(-D)\to(\mathcal O_{\bar X}/\mathcal O_{\bar X}(-D))^2$ 将 $L/L(-D)$ 映到
$(\mathcal O_X/\mathcal O_X(-D))^2$ 的第一分量。容易看出，每个带适配水平结构的 $FH$ 扩张，都能按上述方法得到，成为 $G_D^{L,H}$ 或 $G_D^{H,L}$ 上的点。

2. 另一方面，$\GL_2(\mathcal O_X/\mathcal O_X(-D))$ 作用于带有沿 $D$ 的水平结构的 $FH$ 层概形（1.4）。因此，对仿射直线 $G_D^{L,H}$ 与 $G_D^{H,L}$ 施加这一群作用，可得到模概形上的其他仿射直线。这些直线的并上的点，对应于水平结构不一定适配的 $FH$ 扩张。

3. 若去掉 $D$ 的次数必须很大的假设，本质上没有变化。可先转到一个更大的除子，再对相应的群作用取商。所得是仿射直线对有限群的商，仍为仿射直线。

4. 若 $E/\bar X$ 是 $FH$ 层，含有满足 $\deg F>\frac12\deg E+1$ 的子丛 $F\subset E$，则容易看出 $\Phi$ 必须把 $F^\sigma$ 映到 $F'$，所以 $E$ 必为 $FH$ 扩张。若整数 $s\geq2$，$M_D^{(v,w)}(s)$ 表示由满足下述条件的向量丛 $E$ 给出的子概形：对每个线子丛 $F\subset E$ 都有
$\deg F\leq\frac12\deg E-s$，则
\[
M_D^{(v,w)}(s-1)-M_D^{(v,w)}(s)
\]
是左、右 $FH$ 扩张的仿射直线的并；其中的子丛遍历所有次数为 $s$、秩为一的 $FH$ 层（分别为定义在 $\F_q$ 上的层）。\footnote{\textbf{译者注（本次校订）：}这里的截断符号与前一句不稳定丛的判据不符。要按不稳定程度分层，应以 $\deg F\leq\tfrac12\deg E+s$ 定义 $M_D^{(v,w)}(s)$，并考察 $M_D^{(v,w)}(s)-M_D^{(v,w)}(s-1)$；当 $s\geq2$ 时，新增的丛才由前面的判据保证是 $FH$ 扩张。原式中的减号筛选的是越来越稳定的丛。“子丛次数为 $s$”还须结合所固定的 $\deg E$ 作相应平移。}

5. 当 $\deg L$（分别为 $\deg H_1$）很小时，仿射直线 $G_D^{L,H}$（分别为 $G_D^{H_1,L_1}$）含有一段非空开部分，其中的丛是稳定的。但是当 $\deg L$（分别为 $\deg H_1$）趋向 $-\infty$ 时，不稳定点会越来越多，且不稳定程度越来越高。
\origsection{2}{模空间上的几何点（The geometric points on the moduli space）}
本节的主要结果将模空间某些开子概形上的点数，与尖点形式空间（space of cusp forms）上的 Hecke 算子（Hecke operators）的迹联系起来。
% 原书第364页
\origsubsection{2.1}{基本公式（The fundamental formula）}
始终假设 $v,w\in X(\bar\F_q)$，并且它们在 $K$ 上诱导不同赋值 $\mathfrak p,\mathfrak q$。考虑概形
$M_D^{(v,w)}\to\Spec(\F_{q^d})$，这里 $\F_{q^d}$ 是 $\mathfrak p,\mathfrak q$ 的剩余域的合成域。若 $E(\psi,\alpha_v,\alpha_w,\Phi)$ 是 $FH$ 扩张，则有两个正合列
\[
0\to L\to E\to H\to0,\qquad0\to H_1\to E\to L_1\to0.
\]
若 $\deg L>\deg H$ 或 $\deg H_1\geq\deg L_1$，称这个扩张为强分解的（strongly decomposed）。不是强分解的 $FH$ 层构成开子概形
$U_D\hookrightarrow M_D^{(v,w)}$。群 $P_D$ 由 $X$ 上带沿 $D$ 的水平结构的线丛构成，它通过张量积（tensorisation）作用于 $U_D$ 与 $M_D^{(v,w)}$，从而在紧支撑上同调（cohomology with compact supports）
$H_!^i(\bar U_D,\bar\Q_\ell)$ 上诱导同一群的作用，其中
$\bar U_D=U_D\times_{\F_{q^d}}\bar\F_q$，$\bar\Q_\ell$ 是 $\Q_\ell$ 的代数闭包（algebraic closure）。令 $\mu\subset\bar\Q_\ell$ 为单位根群（group of roots of unity）。对于任意有限特征（finite character）$\omega:P_D\to\mu$，以 $H_!^i(\bar U_D,\bar\Q_\ell)_\omega$ 表示相应的特征部分（eigenspace）；这一记号按2.3.5中对无限循环群取商后的紧支撑上同调来理解。我们关心几何 Frobenius 的幂 $\operatorname{Fr}^m$（$d\mid m$）在上同调上的迹，也就是各次迹的交错和（alternating sum）：
\[
\operatorname{trace}\operatorname{Fr}^m|H_!^\bullet(\bar U_D,\bar\Q_\ell)_\omega
=\sum_{i=0}^4(-1)^i\operatorname{trace}\operatorname{Fr}^m|H_!^i(\bar U_D,\bar\Q_\ell)_\omega.
\]
\paragraph*{注（Remark）}
概形 $U_D\to\Spec(\F_{q^d})$ 不是有限型的（of finite type）。有 1.5.2 引入的行列式映射 $\det:U_D\to\widetilde P_D$。对于
$\xi\in\widetilde P_D(\bar\F_q)$，纤维 $U_D^{(\xi)}$ 是拟射影概形；若 $a\in P_D$，则张量积定义映射
$U_D^{(\xi)}\to U_D^{(\xi a^2)}$。

若 $P_{D,\omega}$ 是 $\omega$ 的核，则
\[
H_!^\bullet(\bar U_D,\bar\Q_\ell)_\omega
\subset H_!^\bullet(\bar U_D,\bar\Q_\ell)^{P_{D,\omega}}
=\bigoplus_{\xi\in\widetilde P_D/P_D^{2,\omega}}
H_!^\bullet(\bar U_D^{(\xi)},\bar\Q_\ell).
\]
因此，上述迹所含的信息等价于关于 $U_D^{(\xi)}$ 的有理点数（number of rational points）的某些信息。

下面把 $\operatorname{Fr}^m$ 的这个迹，与自守形式空间上某些 Hecke 算子的迹联系起来。

选定嵌入（embedding）$\mu\hookrightarrow\C^*$，于是给定的特征 $\omega$ 也可视为取值于 $\C^*$ 的特征。置
$\mathcal K_{\mathfrak p}=\GL_2(\mathcal O_{\mathfrak p})$、
$\mathcal K_{\mathfrak q}=\GL_2(\mathcal O_{\mathfrak q})$，并令
$\mathcal K'_D\subset\prod_{\mathfrak p'\ne\mathfrak p,\mathfrak q}\GL_2(\mathcal O_{\mathfrak p'})$
为 $D$ 所定义的全同余子群（full congruence subgroup）。若 $Z(A)$ 为 $\GL_2(A)$ 的中心（centre），则有等同
\[
Z(A)\big/\bigl[Z(K)\cdot(Z(A)\cap(\mathcal K_{\mathfrak p}\times\mathcal K_{\mathfrak q}\times\mathcal K'_D))\bigr]
\simeq P_D.
\]
因此也可把 $\omega$ 视为 $Z(A)$ 上的复值特征。引入空间
\[
H_\omega=L^2_{\omega,\mathrm{disc}}
\bigl(\GL_2(K)\backslash\GL_2(A)/(\mathcal K'_D\times\mathcal K_{\mathfrak p}\times\mathcal K_{\mathfrak q})\bigr),
\]
其中 $\omega$ 是上述中心特征（central character），下标 $\mathrm{disc}$ 表示只考虑离散谱（discrete spectrum）[1, §4]。该中心特征在 $\mathfrak p,\mathfrak q$ 两处均非分歧（unramified）。引入 Hecke 算子代数（algebras of Hecke operators）
$\mathcal H_{\mathfrak p,\omega}$、$\mathcal H_{\mathfrak q,\omega}$。

$\mathcal H_{\mathfrak p,\omega}$ 的元素，是函数 $f:\GL_2(K_{\mathfrak p})\to\C$，它们模中心紧支撑（compactly supported modulo the centre），满足
\[
f(xz)=f(x)\omega(z),\qquad z\in Z(K_{\mathfrak p}),\quad x\in\GL_2(K_{\mathfrak p}),
\]
% 原书第365页
并且为 $\mathcal K_{\mathfrak p}$ 双不变（bi-invariant）。每个这样的 $f$ 定义算子
$T_f:H_\omega\to H_\omega$：
\[
T_f(h)(y)=\int_{\GL_2(K_{\mathfrak p})/Z(K_{\mathfrak p})}
h(yx)f(x^{-1})\,dx,
\]
其中 $y\in\GL_2(A)$，而 $\GL_2(K_{\mathfrak p})\hookrightarrow\GL_2(A)$。

在这些代数中，对每个 $n\in\Z$，有如下定义的元素 $\Phi_n^{(\mathfrak p)}$。若
$\eta_{s_1,s_2}\ne0$ 是 $\GL_2(K_{\mathfrak p})$ 上满足
\[
\eta_{s_1,s_2}\left(\begin{pmatrix}t_1&*\\0&t_2\end{pmatrix}k_{\mathfrak p}\right)
=|t_1|_{\mathfrak p}^{s_1}|t_2|_{\mathfrak p}^{s_2}
\left|\frac{t_1}{t_2}\right|_{\mathfrak p}^{1/2},
\qquad k_{\mathfrak p}\in\GL_2(\mathcal O_{\mathfrak p}),
\]
且中心特征为 $\omega$ 的函数，即 $\omega(t)=|t|_{\mathfrak p}^{s_1+s_2}$，则
\[
\int_{Z(K_{\mathfrak p})\backslash\GL_2(K_{\mathfrak p})}
\eta_{s_1,s_2}(yx)\Phi_n^{(\mathfrak p)}(x^{-1})\,dx
=(N\mathfrak p)^{|n|/2}
\bigl(N\mathfrak p^{ns_1}+N\mathfrak p^{ns_2}\bigr)\eta_{s_1,s_2}(y),
\]
这里 $N\mathfrak p$ 为 $\mathfrak p$ 处剩余域的基数（cardinality）。这个定义也可用表示（representations）的语言表述：若
$\rho_{\mathfrak p}:\GL_2(K_{\mathfrak p})\to\End(H)$ 是中心特征为 $\omega$ 的第一类不可约表示（irreducible representation of class one），则\footnote{\textbf{译者注（本次校订）：}下式在抽象表示空间中应以 $\eta$ 为右侧最后的向量；$\eta_{s_1,s_2}$ 是前面诱导模型中的函数。这里的含义是 $\mathcal K_{\mathfrak p}$ 不变直线上 Hecke 算子的标量特征值。}
\[
\rho_{\mathfrak p}(T_{\Phi_n^{(\mathfrak p)}})\eta
=(N\mathfrak p)^{|n|/2}
\bigl(N\mathfrak p^{ns_1}+N\mathfrak p^{ns_2}\bigr)\eta_{s_1,s_2},
\]

其中 $\eta\ne0$ 为 $\mathcal K_{\mathfrak p}$ 不变向量，$s_1,s_2$ 是这个表示的参数（参见 [1, §3, B]）。

有了上述记号，便可陈述主要等式。以 $d_{\mathfrak p}$（分别为 $d_{\mathfrak q}$）表示素位 $\mathfrak p$（分别为 $\mathfrak q$）的次数。算子
\[
\Phi_m=T_{\Phi_{m/d_{\mathfrak p}}^{(\mathfrak p)}}\circ
T_{\Phi_{-m/d_{\mathfrak q}}^{(\mathfrak q)}}:H_\omega\to H_\omega
\]
为迹类算子（of trace class；参见 [1, §2]），而且
\begin{align*}
\operatorname{trace}\operatorname{Fr}^m|H_!^\bullet(\bar U_D,\bar\Q_\ell)_\omega
={}&\operatorname{trace}\Phi_m|H_\omega\\
&-\text{算子 }\Phi_m\text{ 的迹公式中的 Eisenstein 贡献}.
\end{align*}
Eisenstein 贡献是 [1] 中第 6.36、6.37 项之和。我们这里的记号与 [1] 不完全一致，请读者参见 §2.3。
\origsubsection{2.2}{主要等式的推论（Consequences of the main equality）}
讨论主要等式的证明之前，先给出若干应用。它既给出概形 $U_D$ 的若干几何性质，也蕴含尖点形式的 Ramanujan 猜想。

取 $X/\F_q$ 上次数为一的线丛 $\eta$。它生成无限循环群（infinite cyclic group）$\{\eta\}\subset P_D$，且
$P_D/\{\eta\}=P_{0,D}$ 为有限群。于是不难看出
\[
\sum_{\omega\in\widehat P_{0,D}}
\operatorname{trace}\operatorname{Fr}^m|H_!^\bullet(\bar U_D,\bar\Q_\ell)_\omega
=\operatorname{trace}\operatorname{Fr}^m
\left(\bigoplus_{\xi\in\widetilde P_D;\,\deg\xi=0,1}
H_!^\bullet(\bar U_D^{(\xi)},\bar\Q_\ell)\right).
\]
用 Frobenius 在上同调 $H_!^i(\bar U_D^{(\xi)},\bar\Q_\ell)$ 上作用的特征值（eigenvalues）$\alpha_{i,\nu}$ 来表示这个迹，就得到
$\sum_{i=1}^4\sum_{\nu=1}^{b_i}(-1)^i\alpha_{i,\nu}^m$。

下面对 $\Phi_m$ 在各空间 $H_\omega$ 上的迹作类似展开。已知 $H_\omega$ 还可以进一步分解如下：
% 原书第366页
\[
H_\omega=H_{\omega,\mathrm{cusp}}\oplus H_{\omega,1}
=\bigoplus_\rho H_{\omega,\rho}\oplus\bigoplus_\chi\C\cdot\chi.
\]
\footnote{译注：所供扫描本此行在 $H_{\omega,1}$ 与后面的直和之间未清楚显示等号；这里将后半式写为前半式的分解。}
其中 $\rho$ 遍历尖点形式空间
$L^2_{\omega,\mathrm{cusp}}(\GL_2(K)\backslash\GL_2(A))$
内 $\GL_2(A)$ 的有限个不可约表示；$\chi:\GL_2(A)\to\C$ 是经由行列式分解的函数，分解如下：
\[
\GL_2(A)\xrightarrow{\det}I_K\longrightarrow P_D
\xrightarrow{\omega'}\mu\subset\C^*,\qquad(\omega')^2=\omega.
\]
分别计算 $\Phi_m$ 在这些直和分量上的迹。

表示 $\rho$ 在 $\mathfrak p,\mathfrak q$ 处必须为第一类，因此 $\sigma_{\mathfrak p}$ 由表示
\[
\rho_{s_1,s_2}:\begin{pmatrix}t_{1,\mathfrak p}&*\\0&t_{2,\mathfrak p}\end{pmatrix}
\longmapsto|t_{1,\mathfrak p}|_{\mathfrak p}^{s_1}|t_{2,\mathfrak p}|_{\mathfrak p}^{s_2}
|t_{1,\mathfrak p}/t_{2,\mathfrak p}|_{\mathfrak p}^{1/2}
\]
诱导得到。定义 $\xi_{\mathfrak p}(\rho)=q^{s_1/d_{\mathfrak p}}$、
$\xi'_{\mathfrak p}(\rho)=q^{s_2/d_{\mathfrak p}}$ 等，则
\begin{align*}
\operatorname{trace}\Phi_m|H_{\omega,\rho}
={}&(\dim H_{\omega,\rho}^{\mathcal K})q^m
\bigl(\xi_{\mathfrak p}(\rho)^m+\xi'_{\mathfrak p}(\rho)^m\bigr)\\
&\hspace{3em}\cdot\bigl(\xi_{\mathfrak q}(\rho)^m+\xi'_{\mathfrak q}(\rho)^m\bigr).
\end{align*}
这应由 $\Phi_m$ 的定义直接得到。\footnote{\textbf{译者注（本次校订）：}按2.1的定义及 $N\mathfrak p=q^{d_{\mathfrak p}}$，参数应取 $\xi_{\mathfrak p}=q^{s_{1,\mathfrak p}}$、$\xi'_{\mathfrak p}=q^{s_{2,\mathfrak p}}$，而 $\mathfrak q$ 处因使用负指标，应取 $\xi_{\mathfrak q}=q^{-s_{1,\mathfrak q}}$、$\xi'_{\mathfrak q}=q^{-s_{2,\mathfrak q}}$。例如 $(N\mathfrak p)^{(m/d_{\mathfrak p})s_1}=q^{ms_1}$，不应再次除以 $d_{\mathfrak p}$。如此定义，所写的迹公式才与 $\Phi_m$ 一致。}

一维空间 $\C_\chi$ 的贡献也很容易计算。函数 $\chi$ 为
\[
\chi\begin{pmatrix}t_{1,\mathfrak p}&*\\0&t_{2,\mathfrak p}\end{pmatrix}
\longmapsto\chi(t_{1,\mathfrak p}t_{2,\mathfrak p})
=|t_{1,\mathfrak p}|_{\mathfrak p}^{s-1/2}
|t_{2,\mathfrak p}|_{\mathfrak p}^{s+1/2}
|t_{1,\mathfrak p}/t_{2,\mathfrak p}|_{\mathfrak p}^{1/2}.
\]
置 $\eta_{\mathfrak p}=\chi(\pi_{\mathfrak p})^{1/d_{\mathfrak p}}$ 等，由定义得到\footnote{\textbf{译者注（本次校订）：}同一归一化下，下式应读作 $\Phi_m$ 对 $\chi$ 的特征值乘以向量 $\chi$；其中须取 $\eta_{\mathfrak p}=\chi(\pi_{\mathfrak p})^{-1/d_{\mathfrak p}}$、$\eta_{\mathfrak q}=\chi(\pi_{\mathfrak q})^{1/d_{\mathfrak q}}$。这是因为 $\chi(\pi_{\mathfrak p})=(N\mathfrak p)^{-s}$，而正指标算子给出 $q^{ms}(1+q^m)$。当 $m$ 被相应次数整除时，根的选择不影响这些 $m$ 次幂。}
\[
T_{\Phi_m}\chi=
(\eta_{\mathfrak p}^m+\eta_{\mathfrak p}^mq^m)
(\eta_{\mathfrak q}^m+\eta_{\mathfrak q}^mq^m).
\]

回顾 $\eta_{\mathfrak p},\eta_{\mathfrak q}$ 均为单位根。

直接从迹公式可以验证，$\Phi_m$ 的迹中的 Eisenstein 贡献可写成
\[
\sum_{i=0}^S\varepsilon_i\eta_i^m
=\text{算子 }\Phi_m\text{ 的 Selberg 迹公式中的 Eisenstein 贡献},
\]
其中 $\varepsilon_i=\pm1$，$\eta_i$ 为权（weight）$0,1,2$ 的代数整数（algebraic integers），即
$|\eta_i|=q^{\nu_i}$，$\nu_i=0,1,2$（参见 [8, 3.1]）。\footnote{\textbf{译者注（本次校订）：}按照下文采用的权的约定，此处应为 $|\eta_i|=q^{\nu_i/2}$，$\nu_i=0,1,2$。原式 $q^{\nu_i}$ 对应权 $0,2,4$，与后文“Eisenstein 贡献中没有权四项”的论证不符。}

定义 $a(\rho)=\dim H_{\omega,\rho}^{\mathcal K}$，则主要等式为
\begin{align*}
\sum_{i=0}^4\sum_{\nu=1}^{b_i}(-1)^i\alpha_{i,\nu}^m
={}&\sum_\rho a(\rho)q^m
\bigl(\xi_{\mathfrak p}(\rho)^m+\xi'_{\mathfrak p}(\rho)^m\bigr)
\bigl(\xi_{\mathfrak q}(\rho)^m+\xi'_{\mathfrak q}(\rho)^m\bigr)\\
&+\sum_\chi(\eta_{\mathfrak p}^m+\eta_{\mathfrak p}^mq^m)
(\eta_{\mathfrak q}^m+\eta_{\mathfrak q}^mq^m)
+\sum_{i=0}^S\varepsilon_i\eta_i^m.
\end{align*}
下面利用这个等式推出若干结论。先回顾 Frobenius 特征值 $\alpha_{i,\nu}$ 的已知性质。

$\alpha_{i,\nu}$ 是代数整数，且 $|\alpha_{i,\nu}|=q^{j/2}$，其中整数 $j\leq i$。这个数 $j$ 就是 $\alpha_{i,\nu}$ 的权（参见 [9, 14]）。另一方面，有下面的简单引理。
\begin{originaltheorem}{引理 2.2.1（Lemma 2.2.1）}
若 $x_1,\ldots,x_n,x'_1,\ldots,x'_{n'},y_1,\ldots,y_m,y'_1,\ldots,y'_{m'}$ 是复数，且
\[
\sum_{i=1}^nx_i^k-\sum_{i=1}^my_i^k
=\sum_{i=1}^{n'}(x'_i)^k-\sum_{i=1}^{m'}(y'_i)^k,
\qquad k=1,2,\ldots,
\]
对所有 $i,j$ 都有 $x_i\ne y_j$ 且 $x'_i\ne y'_j$，则 $n=n'$、$m=m'$，并且在调整顺序后有
$x_i=x'_i$、$y_i=y'_i$。\footnote{\textbf{译者注（本次校订）：}引理还须假设所列复数均非零。正整数次幂和无法检测零元素，例如在任一侧额外添入一个 $0$，所有等式仍成立而项数改变。这里用于 Frobenius 及 Hecke 参数时，各数非零，故可应用补足后的引理。}
\end{originaltheorem}
% 原书第367页
由于 $\omega$ 是有限特征，已知 $\xi_{\mathfrak p}(\rho)\xi'_{\mathfrak p}(\rho)$ 与
$\xi_{\mathfrak q}(\rho)\xi'_{\mathfrak q}(\rho)$ 均为单位根。此外，由自守形式理论有
$|\xi_{\mathfrak p}(\rho)|<q$、$|\xi_{\mathfrak q}(\rho)|<q$。\footnote{\textbf{译者注（本次校订）：}后面的权比较需要更强的界 $|\xi_{\mathfrak p}|,|\xi'_{\mathfrak p}|,|\xi_{\mathfrak q}|,|\xi'_{\mathfrak q}|<q^{1/2}$。这是球形尖点局部分量的酉性及非一维性给出的严格界，对应于归一化诱导参数 $|\Re s_i|<1/2$。仅用正文所写的 $<q$，尚不能排除尖点项中出现权四项。}这说明来自 $\C\cdot\chi$ 的贡献中的项
$(\eta_{\mathfrak p}\eta_{\mathfrak q})^mq^{2m}$ 是右侧仅有的权四项，其项数为 $2|P_{0,D}|$。具体地说，对每个为平方的 $\omega$，在 $P_D$ 的特征群中，方程 $(\omega')^2=\omega$ 有
$2|P_{0,D}/P_{0,D}^2|$ 个解。

另一方面，左侧权四项来自最高次上同调（top cohomology）：
\[
\sum_{\xi,\,\deg(\xi)=0,1}
\operatorname{trace}\operatorname{Fr}^m|H_!^4(\bar U_D^{(\xi)},\bar\Q_\ell)
=\sum_{\nu=1}^{b_4}\alpha_{4,\nu}^m.
\]
下一节将看到，$\bar U_D^{(\xi)}/\bar\F_q$ 都非空。$\xi$ 的个数同样为 $2|P_D|$。\footnote{\textbf{译者注（本次校订）：}此处应为 $2|P_{0,D}|$。群 $P_D$ 含有次数方向的无限循环部分，而这里仅计数次数为0或1的行列式类，每个次数有 $|P_{0,D}|$ 个。}因而得到
$\bar U_D^{(\xi)}/\bar\F_q$ 不可约（irreducible）。

另一方面，左侧权三项带负号，因为它们必须来自
$H_!^3(\bar U_D^{(\xi)},\bar\Q_\ell)$。但关于 Eisenstein 贡献的说明表明，右侧没有带负号的权三项。因此左侧也不能有这样的项。不难看出，这等价于：对任意光滑紧化（smooth compactification）
$\bar U_D^{(\xi)}\hookrightarrow Y$，都有
\[
H^3(\bar Y,\Q_\ell)=H^1(\bar Y,\Q_\ell)=0.
\]
不存在权三项，随即蕴含 $q^m\xi_{\mathfrak p}(\rho)$ 等项必须为权二项，即 $|\xi(\rho)|=1$，等等。\footnote{\textbf{译者注（本次校订）：}这里实际需要比较的是 $q\,\xi_{\mathfrak p}^{(*)}\xi_{\mathfrak q}^{(*)}$ 的四个乘积，正文漏写了 $\mathfrak q$ 处的因子。令 $r_{\mathfrak p},r_{\mathfrak q}$ 为两处参数的最大绝对值；因每对参数之积为单位根，二者均至少为1。由前述严格界，$q r_{\mathfrak p}r_{\mathfrak q}<q^2$；若它大于 $q$，权的整性迫使它等于 $q^{3/2}$，但权三项已被排除。因此 $r_{\mathfrak p}=r_{\mathfrak q}=1$，得到所需结论。}

总结如下。
\begin{originaltheorem}{定理 2.2.2（Theorem 2.2.2）}
\textup{(1)} 概形 $U_D^{(\xi)}/\F_{q^d}$ 绝对不可约（absolutely irreducible）。

\textup{(2)} 上同调群 $H_!^3(\bar U_D^{(\xi)},\Q_\ell)$ 不含权为 $3$ 的类。
\end{originaltheorem}
此外还得到：

\emph{函数域上 $\GL_2(A)$ 的自守尖点形式满足 Ramanujan 猜想，即对任意出现在 $H_{\omega,\mathrm{cusp}}$ 中的 $\rho$，其局部分量（local components）$\rho_{\mathfrak p},\rho_{\mathfrak q}$ 都属于酉主系列（unitary principal series）。}
\paragraph*{注（Remark）}
Drinfeld 本人用另一个论证来证明定理 2.2.2 的 (2)。这个结论也可由 $U_D/\F_q$ 中存在两族仿射直线推出。
\origsubsection{2.3}{主要等式的证明（The proof of the main equality）}
主要等式的证明与 Shimura 簇（Shimura varieties）理论中的方法相似（参见 [3]、[5]）。先描述模空间的几何点集，以及 Frobenius 和 Picard 群（Picard group）在这个集合上的作用。我们将几何点集分解为若干“轨道”，分别计算各轨道对 Frobenius 迹的贡献，再与迹公式中相应的 Hecke 算子迹的贡献比较。
\paragraph*{2.3.1. 几何点集（The set of geometric points）}
为简单起见，假设 $v,w\in X(\bar\F_q)$ 实际上为有理点，即 $\mathfrak p,\mathfrak q$ 的次数均为一。这一假设简化了记号，证明的基本步骤仍适用于一般情形。
% 原书第368页
此时模空间 $M_D^{(v,w)}$ 定义在 $\F_q$ 上。由于 $v,w$ 已固定，以下略去上标 $(v,w)$。$\bar X$ 上带有沿 $D$ 的水平结构的向量丛，与双陪集（double cosets）
$\GL_2(K_\infty)\backslash\GL_2(A_\infty)/\mathcal K_\infty(D)$ 一一对应（参见 1.2）。因此集合
\[
M(\bar\F_q)=\mathop{\operatorname{proj\,lim}}_{D;\,v,w\notin\operatorname{supp}(D)}M_D(\bar\F_q)
\]
是下面双陪集空间的子集：
\[
M(\bar\F_q)\subset
\GL_2(K_\infty)\backslash\GL_2(A_\infty)/(\GL_2(\mathcal O_v)\times\GL_2(\mathcal O_w)).
\]
\paragraph*{注（Remark）}
记号参见 1.1。仍把 $v,w$ 看作 $K_\infty$ 的赋值，因此 $K_v,K_w$ 是 $K_\infty$ 的完备化，而 $\mathcal O_v,\mathcal O_w$ 为相应整数环。

$\boldsymbol x\in\GL_2(A_\infty)$ 给出 $M(\bar\F_q)$ 的点，当且仅当下式有解：
\[
\boldsymbol x^\sigma=\gamma^{-1}\boldsymbol x\boldsymbol\alpha_v
\widetilde{\boldsymbol\alpha}_w.\tag{*}
\]
这里 $\sigma$ 由 Frobenius 在 $K_\infty$ 常数上的作用诱导，$\gamma\in\GL_2(K_\infty)$，并且
\[
\alpha_v\in\GL_2(\mathcal O_v)\begin{pmatrix}\pi_v&0\\0&1\end{pmatrix}\GL_2(\mathcal O_v),
\qquad
\widetilde\alpha_w\in\GL_2(\mathcal O_w)\begin{pmatrix}\pi_w^{-1}&0\\0&1\end{pmatrix}\GL_2(\mathcal O_w).
\]
$\boldsymbol\alpha_v$（分别为 $\widetilde{\boldsymbol\alpha}_w$）是 $v$ 分量（分别为 $w$ 分量）为 $\alpha_v$（分别为 $\widetilde\alpha_w$）、其余分量为 $1$ 的阿代尔。注意 $\gamma,\alpha_v,\widetilde\alpha_w$ 由 $\boldsymbol x$ 确定。

若 (*) 有两个解
\[
\boldsymbol x^\sigma=\gamma^{-1}\boldsymbol x\boldsymbol\alpha_v\widetilde{\boldsymbol\alpha}_w,
\qquad
(\boldsymbol x')^\sigma=(\gamma')^{-1}\boldsymbol x'\boldsymbol\alpha'_v
\widetilde{\boldsymbol\alpha}'_w,
\]
则 $\boldsymbol x,\boldsymbol x'$ 对应于 $M(\bar\F_q)$ 的同一点，当且仅当
$\boldsymbol x=a^{-1}\boldsymbol x'k_vk_w$，其中
$a\in\GL_2(K_\infty)$，$k_v\in\GL_2(\mathcal O_v)$，$k_w\in\GL_2(\mathcal O_w)$。于是
\[
\gamma'=a\gamma a^{-\sigma},\qquad
\alpha'_v=k_v\alpha_vk_v^{-\sigma},\qquad
\widetilde\alpha'_w=k_w\widetilde\alpha_wk_w^{-\sigma}.
\]
这里利用了 $v,w$ 为有理点的假设。关系 $\gamma'=a\gamma a^{-\sigma}$ 意味着 $\gamma,\gamma'$ 是 $\sigma$ 共轭的（$\sigma$-conjugate）；也称 $\alpha'_v$ 与 $\alpha_v$ 为 $\sigma$-$\GL_2(\mathcal O_v)$ 共轭，$\widetilde\alpha'_w$ 与 $\widetilde\alpha_w$ 为 $\sigma$-$\GL_2(\mathcal O_w)$ 共轭。由此得到：

$M(\bar\F_q)$ 的每个几何点给出以下数据：
\begin{enumerate}[label=(\roman*)]
\item $\GL_2(K_\infty)$ 中的一个 $\sigma$ 共轭类（$\sigma$-conjugacy class）$\{\gamma\}$；
\item $\GL_2(\mathcal O_v)\begin{pmatrix}\pi_v&0\\0&1\end{pmatrix}\GL_2(\mathcal O_v)$ 中的一个
$\sigma$-$\GL_2(\mathcal O_v)$ 共轭类 $\{\alpha_v\}$；
\item $\GL_2(\mathcal O_w)\begin{pmatrix}\pi_w^{-1}&0\\0&1\end{pmatrix}\GL_2(\mathcal O_w)$ 中的一个
$\sigma$-$\GL_2(\mathcal O_w)$ 共轭类 $\{\widetilde\alpha_w\}$。
\end{enumerate}
给定这些数据，存在对应点，当且仅当 $\gamma$ 在 $\GL_2(K_v)$（分别为 $\GL_2(K_w)$）中与 $\alpha_v$（分别为 $\widetilde\alpha_w$）$\sigma$ 共轭，并且在 $v,w$ 之外与 $1$ $\sigma$ 共轭。
% 原书第369页
因此，可按以下三步描述几何点集 $M(\bar\F_q)$。
\begin{enumerate}[label=(\roman*)]
\item 描述 $\GL_2(K_\infty)$ 中的 $\sigma$ 共轭类集；
\item 给定这样的类，描述可能的类 $\alpha_v,\widetilde\alpha_w$；
\item 给定 $\gamma,\alpha_v,\widetilde\alpha_w$，描述属于这组数据的几何点集，即对这些特定数据，求方程
$\boldsymbol x^\sigma=\gamma^{-1}\boldsymbol x\boldsymbol\alpha_v\widetilde{\boldsymbol\alpha}_w$ 的解。
\end{enumerate}
\paragraph*{2.3.2. $\sigma$ 共轭（$\sigma$-conjugacy）}
先讨论 (i)，使用 Saito--Shintani 的范数方法（method of norms），参见 [7]。

若 $\gamma\in\GL_2(K_\infty)$，则存在整数 $n$ 使 $\gamma\in\GL_2(K\cdot\F_{q^n})$。这时定义
$N_n(\gamma)=\gamma\gamma^\sigma\cdots\gamma^{\sigma^{n-1}}=\gamma_0$。不难证明：

(1) $\gamma_0$ 的共轭类定义在 $K$ 上。

(2) 若用 $a\in\GL_2(K\cdot\F_{q^n})$ 作 $\sigma$ 共轭，得到
$a\gamma a^{-\sigma}=\gamma'$，则
$aN_n(\gamma)a^{-1}=N_n(\gamma')$。而且可找到 $a\in\GL_2(K\F_{q^n})$，使得
$N_n(\gamma')=N_n(a\gamma a^{-\sigma})=\gamma_0\in\GL_2(K)$。

还需要下面这个进一步的断言。

(3) 若 $N_n(\gamma),N_n(\gamma')$ 在 $\GL_2(K_\infty)$ 中共轭，则 $\gamma,\gamma'$ 为 $\sigma$ 共轭。

证明 (3) 要用到 Galois 上同调（Galois cohomology）与类域论（class-field theory）。下面的讨论将说明证明所需的主要步骤。

始终假设 $N_n(\gamma)=\gamma_0\in\GL_2(K)$，并进一步假设 $\gamma_0$ 为半单的（semisimple）。这样做不失一般性，因为可将 $n$ 换为 $n\cdot p$，从而把 $\gamma_0$ 换为其 $p$ 次幂。

于是最终有两种可能。若 $\gamma_0$ 位于中心，称此 $\sigma$ 共轭类为中心的（central）。若 $\gamma_0$ 为半单，且对所有 $k\geq1$，幂 $\gamma_0^k$ 都不位于中心，则称此 $\sigma$ 共轭类为非中心的（noncentral）。若 $\gamma_0$ 生成一个域，称该类为非分解的（nondecomposed）；若 $\gamma_0$ 位于分裂环面（split torus）中，则称为分解的（decomposed）。

先考虑中心共轭类。若 $\gamma\in\GL_2(K\F_{q^n})$ 且 $N_n(\gamma)=\gamma_0$ 为中心元素，那么 $\gamma$ 在
$\operatorname{PGL}_2(K\F_{q^n})$ 中的像 $\bar\gamma$ 定义一个
$Z^1(K\F_{q^n}/K,\operatorname{PGL}_2(K\F_{q^n}))$ 中的 $1$ 余循环（$1$-cocycle）。由 Hilbert 第九十定理（Hilbert's Theorem 90），$\gamma,\gamma'$ 为 $\sigma$ 共轭，当且仅当
$\bar\gamma,\bar\gamma'$ 在
$H^1(K\cdot\F_{q^n}/K,\operatorname{PGL}_2(K\F_{q^n}))$ 中定义同一个上同调类。因此，中心 $\sigma$ 共轭类与 $K$ 的 Brauer 群（Brauer group）中的 $2$ 挠元素（elements of exponent $2$）一一对应。\footnote{\textbf{译者注（本次校订）：}这一说法还须固定相应的中心范数数据。若讨论任意中心 $\sigma$ 共轭类，投影到 $\operatorname{PGL}_2$ 会丢失标量信息：$\gamma=aI_2$ 与 $I_2$ 都给出平凡 Brauer 类，但非恒定的 $a\in K^*$ 不能被 $\sigma$ 共轭消去，因为在任一有限常数扩张上取范数会给出 $a^nI_2\ne I_2$。后文 shtuka 的局部条件对范数作了额外限制。}

对于非中心 $\sigma$ 共轭类，这里只记录后文所需的事实：若 $\gamma_0=N_n(\gamma)$ 非中心，则或者
$\gamma_0\in E^\times\hookrightarrow\GL_2(K)$，其中 $E/K$ 为二次扩张（quadratic extension）；或者
$\gamma_0\in T(K)$，其中 $T/K$ 为分裂环面。在两种情形，$\gamma$ 都必须位于 $\gamma_0$ 的中心化子（centralizer）中，即
$\gamma\in E_\infty^\times$ 或 $\gamma\in T(K_\infty)$。
\paragraph*{2.3.3. $\sigma$-$\GL_2(\mathcal O_v)$ 共轭}
下面讨论 (ii)。设给定 $\gamma\in\GL_2(K\F_{q^n})$；问题是，能否解方程
\[
x_v\alpha_vx_v^{-\sigma}=\gamma,\qquad x_v\in\GL_2(K_v),\quad
\alpha_v\in\GL_2(\mathcal O_v)
\begin{pmatrix}\pi_v&0\\0&1\end{pmatrix}\GL_2(\mathcal O_v),
\]
以及解中可能出现哪些 $\alpha_v$ 的 $\sigma$-$\GL_2(\mathcal O_v)$ 共轭类。

第一种情形，假设
\[
N_n(\gamma)=\gamma_0=\varepsilon_{\mathfrak p}\cdot z\in K^\times,
\]
% 原书第370页
其中 $z$ 是中心元素，即 $z\in K^\times$，而 $\varepsilon$ 是 $K_{\mathfrak p}^\times$ 中的单位。可以用 $u\in K_v^\times$ 将 $\gamma$ 改为 $u_v\gamma u_v^{-\sigma}$，使
\[
N_n(u_v\gamma u_v^{-\sigma})
=u_vu_v^{-\sigma^n}N_n(\gamma)
=u_vu_v^{-\sigma^n}\varepsilon_{\mathfrak p}z;
\]
而可选取 $u$ 使 $uu^{-\sigma^n}\varepsilon=1$。

于是得到
$x_vu\alpha_vu^{-\sigma}x_v^{-\sigma}=u\gamma u^{-\sigma}$，并可把
$\alpha_v$ 换成 $u\alpha_vu^{-\sigma}$，因为这不改变其 $\sigma$-$\GL_2(\mathcal O_v)$ 共轭类。从而
\[
N_n(x_v\alpha_vx_v^{-\sigma})
=x_vN(\alpha_v)x_v^{-\sigma^n}=z\in K^\times.
\]
转到射影线性群（projective linear group）后，得到
$\bar x_vN(\bar\alpha_v)\bar x_v^{-\sigma^n}=1$，即
\[
N(\bar\alpha_v)=\bar x_v^{\sigma^n}\bar x_v^{-1}.
\]
若将 $n$ 换为满足所需整除条件的充分大的倍数，则
$\bar x_v^{\sigma^n}\bar x_v^{-1}$ 接近 $1$。因此可找到
$\bar k_v\in\operatorname{PGL}_2(\mathcal O_v)$，使
$\bar k_v^{\sigma^n}\bar k_v^{-1}=\bar x_v^{\sigma^n}\bar x_v^{-1}$。再次用
$k_v\alpha_vk_v^{-\sigma}$ 替换 $\alpha_v$ 后，在原有的
$\sigma$-$\GL_2(\mathcal O_v)$ 共轭类内，可假设
$N_n(\alpha_v)=z_v\in K_v^\times$。容易看出，此时 $2\mid n$，且
\[
N_n(\alpha_v)=\begin{pmatrix}\pi_v^{n/2}&0\\0&\pi_v^{n/2}\end{pmatrix}.
\]
由于已知
$H^1(K\F_{q^n}/K,\operatorname{PGL}_2(\mathcal O(\F_{q^n})))=0$，这种情形的 $\alpha_v$ 与元素
\[
\begin{pmatrix}0&\pi_v\\1&0\end{pmatrix}
\]
为 $\sigma$ 共轭。

下面确定哪些 $\{\gamma\}$ 满足条件
$N_n(\gamma)=\gamma_0=\varepsilon_{\mathfrak p}\cdot z_{\mathfrak p}$？首先，若 $\{\gamma\}$ 为中心类，或者它为非分解类且 $\gamma_0$ 生成的域在 $\mathfrak p$ 处不分裂，那么显然满足条件。在后一情形，只须选 $n$ 为偶数并考虑赋值。称这样的 $\{\gamma\}$ 在 $\mathfrak p$ 处非分解。

设 $\gamma_0$ 非中心且在 $\mathfrak p$ 处分裂。由于 $\gamma$ 必须位于 $\gamma_0$ 的中心化子中，显然 $\gamma_0$ 在 $\GL_2(K_{\mathfrak p})$ 中共轭于
\[
\begin{pmatrix}\pi_v^n&0\\0&1\end{pmatrix},
\]
从而 $\alpha_v$ 必须 $\sigma$-$\GL_2(\mathcal O_v)$ 共轭于
\[
\begin{pmatrix}\pi_v^n&0\\0&1\end{pmatrix}.
\]
\footnote{译注：所供扫描本在以上两个矩阵中都显示指数 $n$；紧随其后的共轭类代表则为指数 $1$，这里保留所见形式。}
称 $\{\gamma\}$ 在 $v$ 处分解。因此，$\alpha_v$ 只有两个可能的
$\sigma$-$\GL_2(\mathcal O_v)$ 共轭类，即
\[
\left\{\begin{pmatrix}\pi_v&0\\0&1\end{pmatrix}\right\},\qquad
\left\{\begin{pmatrix}0&\pi_v\\1&0\end{pmatrix}\right\},
\]
而 $\gamma$ 在 $v$ 处必须与其中之一 $\sigma$ 共轭。
\paragraph*{2.3.4. 轨道（The orbits）}
设给定整体 $\sigma$ 共轭类（global $\sigma$-conjugacy class）$\{\gamma\}$、
$\alpha_v,\widetilde\alpha_w$，及解
$\boldsymbol x_0^\sigma=\gamma^{-1}\boldsymbol x_0\boldsymbol\alpha_v\widetilde{\boldsymbol\alpha}_w$。
利用逼近（approximation），可假设 $x_{0v}\in\GL_2(\mathcal O_v)$、
$x_{0w}\in\GL_2(\mathcal O_w)$。那么具有相同
$\gamma,\alpha_v,\widetilde\alpha_w$ 的其他任意解都形如
$\boldsymbol x=\boldsymbol x_0\boldsymbol y$，其中 $\boldsymbol y$ 满足：
% 原书第371页

(i) 若写成 $\boldsymbol y=\boldsymbol y'y_vy_w$，其中 $\boldsymbol y'$ 在 $v,w$ 处的分量为 $1$，则
$\boldsymbol y'\in\GL_2(A_{\mathfrak p,\mathfrak q})$，即 $\boldsymbol y'$ 在 $\mathfrak p,\mathfrak q$ 之外为有理的；

(ii) $y_v\alpha_vy_v^{-\sigma}=\alpha_v$，
$y_w\widetilde\alpha_wy_w^{-\sigma}=\widetilde\alpha_w$。

定义 $\sigma$ 中心化子（$\sigma$-centralizers）
\begin{align*}
Z_\sigma(\alpha_v)&=\{y_v\in\GL_2(\mathcal O_v)\mid y_v\alpha_vy_v^{-\sigma}=\alpha_v\},\\
Z_\sigma(\widetilde\alpha_w)&=\{y_w\in\GL_2(\mathcal O_w)\mid
 y_w\widetilde\alpha_wy_w^{-\sigma}=\widetilde\alpha_w\}.
\end{align*}
\footnote{译注：以上两处集合的环境群在扫描本中写为 $\GL_2(\mathcal O_v)$、$\GL_2(\mathcal O_w)$；后文又对 $Z_\sigma$ 与这些整数点群取交，并将 $Z_\sigma$ 识别为更大的群。两处定义的环境群疑应为相应的 $\GL_2(K_v)$、$\GL_2(K_w)$。}
方程 $\boldsymbol x^\sigma=\gamma^{-1}\boldsymbol x\boldsymbol\alpha_v\widetilde{\boldsymbol\alpha}_w$
的解，在群
$\GL_2(A_{\mathfrak p,\mathfrak q})\times Z_\sigma(\alpha_v)\times Z_\sigma(\widetilde\alpha_w)$
的作用下形成一个轨道。下面计算稳定子（stabilizer）。首先注意，在非分解情形，
\[
Z_\sigma\left(\begin{pmatrix}0&\pi_v\\1&0\end{pmatrix}\right)=D(K_v)^\times,
\]
其中 $D(K_v)^\times$ 为 $v$ 处四元数除代数（quaternion division algebra）的乘法群，\footnote{\textbf{译者注（本次校订）：}这里的除代数应定义在 $K_{\mathfrak p}$ 上，其乘法群应记为 $D(K_{\mathfrak p})^\times$，与2.3.6中的记号一致。1.1中的 $K_v$ 是常数扩张后的完备域；上述 $\sigma$ 中心化子通过下降给出 $K_{\mathfrak p}$ 上的除代数，并在 $K_v$ 上分裂为矩阵代数。}
$D(K_v)^\times\cap\GL_2(\mathcal O_v)=D^{(1)}(K_v)$ 是既约范数（reduced norm）的绝对值为一的元素所构成的极大紧子群（maximal compact subgroup）。于是
\[
\bar D_v=D(K_v)^\times/D^{(1)}(K_v)=\Z.
\]
在分解情形，
\[
Z_\sigma\left(\begin{pmatrix}\pi_v&0\\0&1\end{pmatrix}\right)
=K_{\mathfrak p}^\times\times K_{\mathfrak p}^\times,
\]
且
\[
Z_\sigma\left(\begin{pmatrix}\pi_v&0\\0&1\end{pmatrix}\right)
\cap\GL_2(\mathcal O_v)=U_{\mathfrak p}\times U_{\mathfrak p},
\]
其中 $U_{\mathfrak p}$ 是单位群（group of units）。这时
\[
\bar D_v=
\frac{Z_\sigma\left(\begin{smallmatrix}\pi_v&0\\0&1\end{smallmatrix}\right)}
{Z_\sigma\left(\begin{smallmatrix}\pi_v&0\\0&1\end{smallmatrix}\right)\cap\GL_2(\mathcal O_v)}
=\Z\times\Z.
\]
由于 $Z_\sigma(\alpha_v)\cap\GL_2(\mathcal O_v)$ 显然平凡地作用于模空间的点，所讨论的点集是
$\GL_2(A_{\mathfrak p,\mathfrak q})\times\bar D_v\times\bar D_w$ 下的一个轨道。$\gamma$ 的 $\sigma$ 中心化子
\[
Z_\sigma(\gamma)=\{a\in\GL_2(K_\infty)\mid a\gamma a^{-\sigma}=\gamma\}
\]
经 $a\mapsto\boldsymbol x_0^{-1}a\boldsymbol x_0$ 映入
$\GL_2(A_{\mathfrak p,\mathfrak q})\times\bar D_v\times\bar D_w$；容易看出，它就是所求的稳定子。因此轨道为
\[
Z_\sigma(\gamma)\backslash
\bigl(\GL_2(A_{\mathfrak p,\mathfrak q})\times\bar D_v\times\bar D_w\bigr).
\]
这样的轨道是 $M(\bar\F_q)$ 的子集，在 Frobenius 的作用下不变。为说明这一点，从
\[
\boldsymbol x^\sigma=\boldsymbol x_0^\sigma\boldsymbol y^\sigma
=\gamma^{-1}\boldsymbol x_0\boldsymbol y\boldsymbol\alpha_v
\widetilde{\boldsymbol\alpha}_w
\]
出发。可以选择 $\alpha_v=\alpha_v^\sigma$、
$\widetilde\alpha_w=\widetilde\alpha_w^\sigma$。于是显然，对于
$\boldsymbol x'=\gamma\boldsymbol x^\sigma$，有
$\boldsymbol x'=\boldsymbol x_0\boldsymbol y\boldsymbol\alpha_v\widetilde{\boldsymbol\alpha}_w$。
这意味着 Frobenius 在上述轨道上的作用，就是右乘
$\alpha_v\widetilde\alpha_w\in\bar D_v\times\bar D_w$。
% 原书第372页
回顾
\[
\begin{aligned}
\alpha_v&=1\in\Z=\bar D_v&&\text{（非分解情形）},\\
\alpha_v&=(1,0)\in\Z\times\Z=D_v&&\text{（分解情形）}.
\end{aligned}
\]
沿 $D$ 带水平结构的有理线丛群作用于 $M_D$（参见 2.1），因而得到
$P=\operatorname{proj\,lim}_DP_D$（$\mathfrak p,\mathfrak q\notin\operatorname{supp}(D)$）在 $M$ 上的作用。注意
\[
P\xrightarrow{\sim}I_K/\bigl[(U_{\mathfrak p}\times U_{\mathfrak q})\cdot K^\times\bigr],
\]
其中 $I_K$ 是伊代尔群（idele group），$U_{\mathfrak p}$（分别为 $U_{\mathfrak q}$）是 $\mathfrak p$（分别为 $\mathfrak q$）处的单位群。由于
$I_K$ 同构于中心 $Z(A)\subset\GL_2(A)$，上述作用在各轨道上由
\[
\boldsymbol x\longmapsto\boldsymbol xz,\qquad z\in Z(A),\quad
\boldsymbol x\in Z_\sigma(\gamma)\backslash
(\GL_2(A_{\mathfrak p,\mathfrak q})\times\bar D_v\times\bar D_w)
\]
给出。

上述讨论说明，$\sigma$ 共轭类 $\{\gamma\}$ 决定其余数据，也决定轨道。还需确定哪些 $\{\gamma\}$ 给出非空轨道，也就是哪些类使方程 (*) 有解。

假设类 $\{\gamma\}$ 为中心类。由
$\boldsymbol x^\sigma=\gamma^{-1}\boldsymbol x\boldsymbol\alpha_v\widetilde{\boldsymbol\alpha}_w$
可知，$\gamma$ 在所有异于 $\mathfrak p,\mathfrak q$ 的素位处都是余边界（boundary）。所以，$\gamma$ 定义的四元数代数，更确切地说中心代数（central algebra），必须在 $\mathfrak p,\mathfrak q$ 处不分裂，而在其他各处都分裂。另一方面，已知恰好存在一个中心 $\sigma$ 共轭类 $\{\gamma\}$，对应于在 $\mathfrak p,\mathfrak q$ 处不分裂的唯一四元数代数。对于这个特定的 $\{\gamma\}$，可取
\[
\alpha_v=\begin{pmatrix}0&\pi_v\\1&0\end{pmatrix},\qquad
\widetilde\alpha_w=\begin{pmatrix}0&\pi_w^{-1}\\1&0\end{pmatrix}
\]
来解 (*)，由此恰好得到一个轨道，称为超奇异轨道（supersingular orbit）。

若 $\gamma$ 非中心，则它决定一个二次域扩张 $E/K$ 或一个分裂环面 $T/K$。始终假设 $E/K$（分别为 $T/K$）以如下方式嵌入 $\GL_2$：
\[
\operatorname{units}(E_{\mathfrak p})\hookrightarrow\GL_2(\mathcal O_{\mathfrak p}),\qquad
\operatorname{units}(E_{\mathfrak q})\hookrightarrow\GL_2(\mathcal O_{\mathfrak q}),
\]
分别地，
\[
\begin{aligned}
\operatorname{units}(T(K_{\mathfrak p}))&=U_{\mathfrak p}\times U_{\mathfrak p}
\hookrightarrow\GL_2(\mathcal O_{\mathfrak p}),\\
\operatorname{units}(T(K_{\mathfrak q}))&=U_{\mathfrak q}\times U_{\mathfrak q}
\hookrightarrow\GL_2(\mathcal O_{\mathfrak q}).
\end{aligned}
\]
由 (*)，当 $n$ 足够大时，有
$\operatorname{div}(N_n(\gamma))=n\mathfrak p-n\mathfrak q$。这蕴含：若 $E/K$ 在 $\mathfrak p,\mathfrak q$ 处都不分裂，则
\[
N_n(\gamma)=a\cdot\varepsilon,\qquad a\in K^\times,\quad
\varepsilon\in E^\times\text{ 是整体单位（global unit）}.
\]
但这时 $\{\gamma\}$ 显然是中心类。因此，在 $\mathfrak p,\mathfrak q$ 两处都不分裂的域扩张，不可能对非超奇异轨道作出贡献。
% 原书第373页
另一方面，将看到，在两处中的至少一处分裂的任意扩张，都将通过一个或两个轨道对点集作出贡献。例如，设 $\mathfrak p$ 在 $E$ 中分解为
$\mathfrak p=\mathfrak p_1\mathfrak p_2$。可找到 $\gamma\in E_\infty$，使其除子形如
\[
\operatorname{div}(\gamma)=-\mathfrak q'+\mathfrak p_1-\mathfrak a+\mathfrak a^\sigma,
\]
这里 $\mathfrak q'$ 是 $E_\infty$ 中位于 $\mathfrak q$ 之上的素位，$\mathfrak a$ 是支撑不含 $\mathfrak p,\mathfrak q$ 的正除子。这样的 $\mathfrak a$ 的存在性，又来自将 Lang 定理应用于 $E$ 的 Jacobian 簇（Jacobian）。

$\gamma$ 的范数的除子为
$\operatorname{div}(N_n(\gamma))=n\mathfrak p_1-n\mathfrak q_1$，其中若 $\mathfrak q$ 不分裂，则 $\mathfrak q_1=\mathfrak q$；若 $\mathfrak q$ 在 $E$ 中分裂，则 $\mathfrak q_1$ 是 $E$ 中位于 $\mathfrak q$ 之上的一个除子。由于 $E/K$ 的非平凡自同构可以实现为 $\GL_2(K)$ 的内自同构（inner automorphism），因此：

每个在 $\mathfrak p$ 或 $\mathfrak q$ 处分裂的域扩张 $E/K$，都对轨道集作出贡献。若恰在一处分裂，它给出一个轨道；若在两处都分裂，则给出两个轨道。
\paragraph*{2.3.5. $\operatorname{Fr}^m$ 的迹的计算（The computation of the trace）}
回顾，我们要计算

\[
\operatorname{trace}\operatorname{Fr}^m|H_!^\bullet(\bar U_D,\Q_\ell)_\omega.
\]
按如下方式进行：特征 $\omega$ 的核中含有无限阶元素 $\lambda\in P_D$，对无限循环群 $\{\lambda\}$ 的作用取商，得到
\[
U_D/\{\lambda\}=\coprod_{\xi;\,\xi\in P_D/\{\lambda^2\}}
U_D^{(\xi)}=Z.
\]
概形 $Z/\F_q$ 为拟射影的（参见 2.2）。在 $Z/\F_q$ 上，有由与线丛作张量积所诱导的有限群
$P_D/\{\lambda\}$ 的作用，并且
\[
H_!^\bullet(\bar U_D,\Q_\ell)_\omega
=H_!^\bullet(\bar Z,\Q_\ell)_\omega
\subset H_!^\bullet(\bar Z,\Q_\ell).
\]
因此，只要能计算 $\operatorname{Fr}^m\circ\delta$ 在
$H_!^\bullet(\bar Z,\Q_\ell)$ 上的迹，其中 $\delta$ 遍历 $P_D/\{\lambda\}$，就得到所需信息。为计算这些迹，必须计算
$\operatorname{Fr}^m\circ\delta$ 在集合 $Z(\bar\F_q)$ 上的不动点数（number of fixed points）。这些不动点可以在前述各个轨道中分别计数。为证明主要等式，将把这些贡献与 Selberg 迹公式中的对应贡献作比较。

回顾 2.3.4 的轨道描述：它是集合
\[
Z_\sigma(\gamma)\cdot\{\lambda\}\backslash
(\GL_2(A_{\mathfrak p,\mathfrak q})\times\bar D_v\times\bar D_w).
\]
元素 $\delta$ 通过右乘伊代尔 $\delta\in Z(A)$ 作用，而
$\operatorname{Fr}^m$ 通过右乘 $(\alpha_v\widetilde\alpha_w)^m$ 作用。令
$\mathcal K'_D\subset\GL_2(A_{\mathfrak p,\mathfrak q})$ 为 $D$ 所定义的开紧子群（open compact subgroup）。定义函数
$\Psi_{D,m,\omega}$ 为下列集合的特征函数（characteristic function）：
\[
\bigcup_{\nu\in\Z}\mathcal K'_D\cdot\lambda^\nu
(\alpha_v\widetilde\alpha_w)^m\cdot\delta\cdot Z(K),
\]
其中 $Z(K)$ 当然是 $\GL_2(K)$ 的中心。直接按共轭类计数，可将轨道上的不动点数写为
\[
\sum_{\substack{a\text{ 为 }Z_\sigma(\gamma)/Z(K)\\\text{ 中的共轭类}}}
\int_{Z_a(K)\backslash(\GL_2(A_{\mathfrak p,\mathfrak q})\times\bar D_v\times\bar D_w)}
\Psi_{D,m,\delta}(x^{-1}ax)\,dx,
\]
其中 $Z_a(K)$ 为元素 $a\in\GL_2(K)$ 在 $\GL_2(K)$ 中的中心化子。Haar 测度（Haar measure）$dx$ 应使开紧子群
$\mathcal K'_D\times\{0\}\times\{0\}$ 的体积为一。
% 原书第374页
由上述不动点数公式，可进一步得到
$\operatorname{Fr}^m$ 在 $H_!^\bullet(\bar U_D,\Q_\ell)_\omega$ 上的迹公式。

定义新函数
$\Psi_{D,m,\omega}:\GL_2(A_{\mathfrak p,\mathfrak q})\times\bar D_v\times\bar D_w\to\C$：
\[
\Psi_{D,m,\omega}(x)=
\begin{cases}
\omega(z),&x=k'_D\cdot z\cdot(\alpha_v\widetilde\alpha_w)^m,
\quad z\in Z(A),\ k'_D\in\mathcal K'_D,\\
0,&\text{其他情形}.
\end{cases}
\]
于是容易看出，这个轨道对 $\operatorname{Fr}^m$ 的迹的贡献由下式给出：
\[
\sum_{\substack{a\text{ 为 }Z_\sigma(\gamma)/Z(K)\\\text{ 中的共轭类}}}
\operatorname{vol}\bigl(Z_a(A)/(Z_a(K)\cdot Z(A))\bigr)\cdot I_{D,m,\omega}(a),
\]
其中
\[
I_{D,m,\omega}(a)=
\int_{Z_a(A)\backslash(\GL_2(A_{\mathfrak p,\mathfrak q})\times\bar D_v\times\bar D_w)}
\Psi_{D,m,\omega}(x^{-1}ax)\,d\bar x.
\]
$d\bar x$ 是上述商上的商测度（quotient measure），使 $dt\,d\bar x=dx$，且
$\operatorname{vol}_{dt}(Z_a(A)\cap(\mathcal K'_D\times\{0\}\times\{0\}))=1$。

上述计算适用于非分解轨道，其所有点均属于 $U(\bar\F_q)$。分解轨道需要单独处理，详见2.3.7。

由于各函数 $\Psi_{D,m,\omega}$ 是局部函数的乘积，积分 $I_{D,m,\omega}(a)$ 也是局部积分的乘积。得到
\begin{align*}
I_{D,m,\omega}(a)
&=\left(\int_{Z_a(A_{\mathfrak p,\mathfrak q})\backslash\GL_2(A_{\mathfrak p,\mathfrak q})}
\chi'_D((x')^{-1}ax')\,d\bar x'\right)\cdot\delta(a,m)\\
&=I'_{D,\omega}(a)\cdot\delta(m,a),
\end{align*}
其中 $\delta(m,a)=\delta_{\mathfrak p}(m,a)\delta_{\mathfrak q}(m,a)$，且
\[
\delta_{\mathfrak p}(m,a)=
\int_{Z_a(K_{\mathfrak p})\backslash Z_\sigma(\alpha_v)}
\Psi_{D,m,\omega,\mathfrak p}(x_{\mathfrak p}^{-1}ax_{\mathfrak p})\,d\bar x_{\mathfrak p}.
\]
仍选取商测度 $d\bar x_{\mathfrak p}=dx_{\mathfrak p}/dt_{\mathfrak p}$，使
\[
\operatorname{vol}_{dt_{\mathfrak p}}(Z_a(K_{\mathfrak p})\cap\GL_2(\mathcal O_{\mathfrak p}))
=\operatorname{vol}_{dx_{\mathfrak p}}(Z_\sigma(\alpha_v)\cap\GL_2(\mathcal O_v))=1.
\]
现在可以与 Selberg 迹公式作比较。
\paragraph*{2.3.6. 与 Selberg 迹公式的比较（Comparison to the Selberg trace formula）}
目前的情形已由 R. Langlands 在若干篇文章中研究过（参见 [3]、[5]）。Hilbert 空间（Hilbert space）
$L^2_{\omega,\mathrm{disc}}(\GL_2(K)\backslash\GL_2(A)/\mathcal K_D)$
上算子 $\Phi_m$ 的迹，将用 Selberg 迹公式来计算。为简单起见，假设特征 $\omega$ 为平凡特征（trivial character）。这会大幅简化记号，以下略去所有下标 $\omega$。

迹公式把迹表示为 $\GL_2(K)/K$ 中不同共轭类的贡献之和，再加上一些额外的项，称为 Eisenstein 项（Eisenstein terms；参见 [1, §6]）。在证明主要等式时，当然可以略去 Eisenstein 项。

先考虑迹公式中来自单位元素（identity element）以及各向异性半单元素（anisotropic semisimple elements）的项。每个这样的半单元素都定义二次扩张 $E/K$。比较结果如下。

(i) $\operatorname{Fr}^m$ 的迹的超奇异贡献 $=$ 单位元素的贡献 $+$ 所有在 $\mathfrak p,\mathfrak q$ 处都不分裂的二次扩张的贡献。（提醒：这些扩张不给出轨道！）
% 原书第375页
(ii) 在 $\mathfrak p$ 或 $\mathfrak q$ 处分裂的扩张对
$\operatorname{Fr}^m$ 的迹的贡献之和 $=$ 同样的扩张对 Selberg 迹的贡献。

先看 (i)。超奇异轨道来自四元数代数 $D/K$，它恰在 $\mathfrak p,\mathfrak q$ 两个素位处不分裂。由 2.3.4 的讨论，
$Z_\sigma(\gamma)=D(K)^*$，$Z_\sigma(\alpha_v)=D(K_{\mathfrak p})^*$，
$Z_\sigma(\widetilde\alpha_w)=D(K_{\mathfrak q})^*$。有
\[
\bar D_v=\Z=D(K_{\mathfrak p})^*/D^{(1)}(K_{\mathfrak p}),\qquad
\bar D_w=\Z=D(K_{\mathfrak q})^*/D^{(1)}(K_{\mathfrak q}),
\]
其中 $D^{(1)}$ 表示既约范数的绝对值应为一。轨道为
\[
D(K)\cdot D^{(1)}(K_{\mathfrak p})\cdot D^{(1)}(K_{\mathfrak q})\backslash D(A)
=D(K)\backslash(\GL_2(A_{\mathfrak p,\mathfrak q})\times\Z\times\Z).
\]
$D(K)^*/K^*$ 中的共轭类，由单位元素以及各 $E^*/K^*$ 中的元素模共轭给出，这里 $E/K$ 遍历在 $\mathfrak p,\mathfrak q$ 两处都不分裂的二次扩张。因此，对 $\operatorname{Fr}^m$ 的迹的贡献为
\[
\operatorname{vol}(D(K)^*\cdot Z(A)\backslash D(A)^*)\cdot\Psi_{D,m}(e)
\]
加上对所有此类二次扩张的求和，其中 $E/K$ 的贡献为
\[
\frac1{|\operatorname{AUT}(E)|}
\sum_{a\in E^*/K^*}\operatorname{vol}(I_E/(I_KE^*))\cdot I'_D(a)\cdot\delta(a,m).
\]
Selberg 迹公式中的贡献非常类似：中心类的贡献为
\[
\operatorname{vol}(\GL_2(K)\cdot Z(A)\backslash\GL_2(A))\cdot\Phi_m(e),
\]
而每个扩张 $E/K$ 给出
\[
\frac1{|\operatorname{AUT}(E)|}
\sum_{a\in E^*/K^*}\operatorname{vol}(I_E/(I_KE^*))\cdot I'_D(a)\cdot\Theta(a,\Phi_m),
\]
其中
$\Theta(a,\Phi_m)=\theta_{\mathfrak p}(a,\Phi_m^{(\mathfrak p)})
\theta_{\mathfrak q}(a,\Phi_m^{(\mathfrak q)})$，且
\[
\theta_{\mathfrak p}(a,\Phi_m^{(\mathfrak p)})
=\int_{Z_a(K_{\mathfrak p})^*\backslash\GL_2(K_{\mathfrak p})}
\Phi_m^{(\mathfrak p)}(x_{\mathfrak p}^{-1}ax_{\mathfrak p})\,d\bar x_{\mathfrak p}.
\]
关键是能够比较因子 $\Psi_{D,m}(e)$ 与 $\Phi_m(e)$，以及
$\delta(a,m)$ 与 $\Theta(a,\Phi_m)$。

仍先看中心类。注意
\[
\operatorname{vol}(D(K)\cdot Z(A)\backslash D(A)^*)
=(q-1)^2\operatorname{vol}(\GL_2(K)\cdot Z(A)\backslash\GL_2(A)).
\]
这个等式来自除代数的射影群与射影线性群具有相同的 Tamagawa 数（Tamagawa numbers）。因子 $(q-1)^2$ 来自测度的正规化：我们使
$\mathcal K'_D\times D^{(1)}(K_{\mathfrak p})\times D^{(1)}(K_{\mathfrak q})$
和 $\mathcal K'_D\times\GL_2(\mathcal O_{\mathfrak p})\times\GL_2(\mathcal O_{\mathfrak q})$
的体积都等于一。

另一方面，显然当 $m$ 为偶数时 $\Psi_{D,m}(e)=1$，为奇数时等于零。我们断言
\[
\Phi_m^{(\mathfrak p)}(e)=
\begin{cases}-(q-1),&m\text{ 为偶数},\\0,&m\text{ 为奇数}.\end{cases}
\]
为验证这个断言，对 $n\geq0$ 引入函数 $f_n\in\mathcal H_{\mathfrak p}$：
% 原书第376页
\[
f_n\left(\begin{pmatrix}\pi_{\mathfrak p}^{n+a}&0\\0&\pi_{\mathfrak p}^a\end{pmatrix}\right)=1,
\]
并在其他 $\GL_2(\mathcal O_{\mathfrak p})$ 双陪集上置 $f_n=0$。于是容易验证
\[
\Phi_n^{(\mathfrak p)}=f_n-(q-1)\sum_{\nu=1}^{[n/2]}f_{n-2\nu}.
\]
此式也可通过对 [4] 引理3.1中的公式作反演得到。\footnote{\textbf{译者注（本次校订）：}此式适用于 $n\geq1$。当 $n=0$ 时，2.1规定的特征值为 $2$，故 $\Phi_0^{(\mathfrak p)}=2f_0$，而不是右边空和所给出的 $f_0$；后文的正整数次 Frobenius 迹计算不涉及这一例外。}它当然蕴含上述断言，且当 $m$ 为偶数时得到
$\Phi_m(e)=(q-1)^2$。由此得到两侧中心类贡献相等。

再取非中心元素 $a\in D_K$，计算
\[
\delta_{\mathfrak p}(a,m)=
\int_{Z_a(K_{\mathfrak p})^*\backslash D(K_{\mathfrak p})^*}
\Psi_{D,m,\mathfrak p}(x_{\mathfrak p}^{-1}ax_{\mathfrak p})\,d\bar x_{\mathfrak p}.
\]
关键在于：共轭类
$\{x_{\mathfrak p}^{-1}ax_{\mathfrak p}\}_{x_{\mathfrak p}\in D(K_{\mathfrak p})^*}$
在极大紧子群 $D^{(1)}(K_{\mathfrak p})$ 的作用下，若 $a$ 生成非分歧扩张，则分成两个轨道；若生成分歧扩张（ramified extension），则只有一个轨道。因此，当 $a$ 生成的扩张在 $\mathfrak p$ 处非分歧时，
\[
\delta_{\mathfrak p}(a,m)=\begin{cases}2,&m\text{ 为偶数},\\0,&m\text{ 为奇数};\end{cases}
\]
若扩张在 $\mathfrak p$ 处分歧，则对所有 $m>0$，有
$\delta_{\mathfrak p}(a,m)=1$。\footnote{\textbf{译者注（本次校订）：}分歧情形还须加上奇偶条件 $m\equiv\operatorname{ord}_{\mathfrak p}\operatorname{Nrd}(a)\pmod2$；不满足时积分为零。因 $\Psi$ 的支撑由标量与 $\alpha_v^m$ 的乘积组成，其约化范数赋值必与 $m$ 同余。下文分歧情形的 $\Theta_{\mathfrak p}$ 同样须附此条件；在 $\GL_2$ 一侧，约化范数就是行列式。}

另一方面，可用 Langlands 在 [4, 3.x] 中的公式，计算轨道积分（orbital integrals）
$\theta_{\mathfrak p}(a,\Phi_m^{(\mathfrak p)})$ 与
$\theta_{\mathfrak q}(a,\Phi_m^{(\mathfrak q)})$。注意，引理 3.5 蕴含，对
$u=\begin{pmatrix}b&1\\0&b\end{pmatrix}$，有
\[
\int_{Z_u(K_{\mathfrak p})\backslash\GL_2(K_{\mathfrak p})}
\Phi_m^{(\mathfrak p)}(x_{\mathfrak p}^{-1}ux_{\mathfrak p})\,d\bar x_{\mathfrak p}=0.
\]
由引理 3.6，若元素 $a$ 在 $\mathfrak p$ 处生成非分歧扩张，则
\[
\theta_{\mathfrak p}(a,\Phi_m^{(\mathfrak p)})
=\begin{cases}2,&m\text{ 为偶数},\\0,&m\text{ 为奇数},\end{cases}
\]
而若 $a$ 生成分歧扩张，则 $\Theta_{\mathfrak p}(a,\Phi_m^{(\mathfrak p)})=1$ [4, 3.7, 3.8]。

因此两侧相应项相等。特征为 $2$ 时还需补充以下核查：当$a$ 生成不可分扩张（inseparable extension）时，计算
$\delta_{\mathfrak p}(a,m)$ 与 $\Theta_{\mathfrak p}(a,\Phi_m^{(\mathfrak p)})$。

下面比较其余非分解轨道。它们来自至少在 $\mathfrak p,\mathfrak q$ 中一处分裂的域扩张。设
$a\in\GL_2(K)$ 生成一个在 $\mathfrak p$ 处分裂的域扩张，需计算
$\delta_{\mathfrak p}(a,m)$ 与 $\theta_{\mathfrak p}(a,\Phi_m^{(\mathfrak p)})$。
由于扩张在 $\mathfrak p$ 处分裂，有 $\bar D_v=\Z\times\Z$，且 $a$ 在 $\bar D_v$ 中的像为
$(\operatorname{ord}_{\mathfrak p'}(a),\operatorname{ord}_{\mathfrak p''}(a))$。
必须在 $\mathfrak p$ 之上的两个素位 $\mathfrak p',\mathfrak p''$ 中选定一个 $\mathfrak p'$，此时
$\operatorname{Fr}^m$ 的作用由 $(m,0)\in D_v$ 给出。这意味着
% 原书第377页
\[
\delta_{\mathfrak p}(a,m)=\delta'_{\mathfrak p}(a,m)
=\begin{cases}1,&\operatorname{ord}_{\mathfrak p'}(a)-\operatorname{ord}_{\mathfrak p''}(a)=m,\\
0,&\text{其他情形}.\end{cases}
\]
另一方面，由 [4, 引理 3.1]，
\[
\Theta_{\mathfrak p}(a,\Phi_m^{(\mathfrak p)})
=\begin{cases}1,&|\operatorname{ord}_{\mathfrak p'}(a)-\operatorname{ord}_{\mathfrak p''}(a)|=m,\\
0,&\text{其他情形}.\end{cases}
\]
假设扩张在 $\mathfrak q$ 处不分裂。也可以选 $\mathfrak p''$ 而非 $\mathfrak p'$，得到同一轨道。因此，若在迹公式左侧使用
$\delta'_{\mathfrak p}(a,m)+\delta''_{\mathfrak p}(a,m)$，得到的是
$\operatorname{Fr}^m$ 的迹的两倍。但是显然
\[
\delta'_{\mathfrak p}(a,m)+\delta''_{\mathfrak p}(a,m)
=\Theta_{\mathfrak p}(a,\Phi_m^{(\mathfrak p)}).
\]
再回顾，Selberg 迹公式一侧有因子 $\frac12$，而 $\operatorname{Fr}^m$ 的公式没有这个因子，就看出给定域对两侧的贡献相等。

若 $a$ 生成的域扩张 $E/K$ 在 $\mathfrak p,\mathfrak q$ 两处都分裂，论证十分相似。2.3.4 已看到，这种情形有两个轨道，对应于在 $\mathfrak p,\mathfrak q$ 之上选取素位的四种可能方式；在 Galois 群作用下共轭的选择给出同一轨道。此后与前面完全相同，即可证明两侧相应的贡献相等。

至此已比较来自中心和椭圆共轭类（elliptic conjugacy classes）的贡献。Lefschetz 不动点公式（Lefschetz fixed point formula）一侧没有幺幂类（unipotent class）的贡献，而这些贡献在 Selberg 迹公式中也消失。前面已经指出，$\Phi_m^{(\mathfrak p)}$ 在幺幂类上的轨道积分为零。由于这在两个素位处都成立，标准论证（参见 [2, 第 523 页]）说明这些项为零。

最后比较分解轨道与分裂半单元素的贡献。这一步需要使用 $U/\F_q$ 的定义，下面单独讨论。
\paragraph*{2.3.7. 分解轨道（The decomposed orbit）}
令 $T/K$ 为 $\GL_2/K$ 中由对角矩阵组成的标准极大环面（standard maximal torus）。以上三角、下三角矩阵组成的两个 Borel 子群，其幺幂根基（unipotent radicals）分别记为
$N^+$（分别为 $N^-$）。选取 $\gamma\in K_\infty^\times$，使
\[
\operatorname{div}(\gamma)=\mathfrak p-\mathfrak q+\mathfrak a-\mathfrak a^\sigma.
\]
若 $\gamma\in K\cdot\F_{q^n}$，则 $\gamma_0=N_n(\gamma)$ 的除子为
$\operatorname{div}(\gamma_0)=n\mathfrak p-n\mathfrak q$。可在 $I_{K_\infty}$ 中找到伊代尔 $x_0$，使
\[
\begin{pmatrix}x_0^\sigma&0\\0&1\end{pmatrix}
=\begin{pmatrix}\gamma^{-1}&0\\0&1\end{pmatrix}
\begin{pmatrix}x_0&0\\0&1\end{pmatrix}
\begin{pmatrix}\pi_v&0\\0&1\end{pmatrix}
\begin{pmatrix}\pi_w^{-1}&0\\0&1\end{pmatrix}.
\]
置
\[
\boldsymbol x_0=\begin{pmatrix}x_0&0\\0&1\end{pmatrix},\qquad
\boldsymbol\gamma=\begin{pmatrix}\gamma&0\\0&1\end{pmatrix}.
\]
$M(\bar\F_q)$ 中属于分解轨道的点，对应于集合
$\boldsymbol x_0\boldsymbol y$ 的元素，其中
$\boldsymbol y\in\GL_2(A_{\mathfrak p,\mathfrak q})\times T(K_{\mathfrak p})\times T(K_{\mathfrak q})$。
轨道为
\[
T(K)\backslash\bigl(\GL_2(A_{\mathfrak p,\mathfrak q})\times(\Z\times\Z)\times(\Z\times\Z)\bigr),
\]
其中 $\Z\times\Z=T(K_{\mathfrak p})/T(\mathcal O_{\mathfrak p})$（分别为
$T(K_{\mathfrak q})/T(\mathcal O_{\mathfrak q})$），而
$\operatorname{Fr}^m$ 通过 $(m,0)\times(-m,0)\in(\Z\times\Z)\times(\Z\times\Z)$ 作用。
现在判断轨道中哪些点属于 $U(\bar\F_q)$。为此，分别作如下分解
% 原书第378页
\[
\boldsymbol y=\boldsymbol t^+(\boldsymbol y)\boldsymbol n^+(\boldsymbol y)
\boldsymbol k'{}^+(\boldsymbol y),
\qquad\text{分别为}\qquad
\boldsymbol y=\boldsymbol t^-(\boldsymbol y)\boldsymbol n^-(\boldsymbol y)
\boldsymbol k'{}^-(\boldsymbol y),
\]
其中 $\boldsymbol t^+(\boldsymbol y),\boldsymbol t^-(\boldsymbol y)\in
T(A_{\mathfrak p,\mathfrak q})\times\Z^4$，
$\boldsymbol n^+(\boldsymbol y)\in N^+(A_{\mathfrak p,\mathfrak q})$
（分别为 $\boldsymbol n^-(\boldsymbol y)\in N^-(A_{\mathfrak p,\mathfrak q})$），
$\boldsymbol k'{}^+(\boldsymbol y),\boldsymbol k'{}^-(\boldsymbol y)\in\mathcal K'$。
这两种写法对应于将相应的 $FH$ 层表示为 $FH$ 扩张的两种方式。为验证这一点，回顾 1.1：这个 $FH$ 层由下列格给出，
\[
\boldsymbol x_0\boldsymbol yM_0
=\boldsymbol x_0\boldsymbol t^+(\boldsymbol y)\boldsymbol n^+(\boldsymbol y)M_0
=\boldsymbol x_0\boldsymbol t^-(\boldsymbol y)\boldsymbol n^-(\boldsymbol y)M_0=M.
\]
每种分解都在 $M$ 中给出一个秩一子层（subsheaf），分别由第一个、第二个基向量生成。称这些子层为 $L$（分别为 $H_1$）。显然 $L$ 是秩一 $FH$ 层，而 $H_1$ 是 $X/\F_q$ 上的可逆层（invertible sheaf）。若写成
\[
\boldsymbol t^+(\boldsymbol y)=\begin{pmatrix}t_1^+&0\\0&t_2^+\end{pmatrix},\qquad
\boldsymbol t^-(\boldsymbol y)=\begin{pmatrix}t_1^-&0\\0&t_2^-\end{pmatrix},
\]
则有（$x_0$ 的次数为零！）
\[
\deg L=-\deg t_1^+=-\log_q|t_1^+|,\qquad\deg H_1=-\deg t_2^-.
\]
得到 $U(\bar\F_q)$ 中的点，当且仅当（参见 2.1）
$|t_1^+|<|t_2^+|$ 且 $|t_2^-|\leq|t_1^-|$。

若 $\pi$ 是 $K$ 的某个素位处的一致化元（uniformizing element），则对所有 $f>0$，有恒等式\footnote{\textbf{译者注（本次校订）：}下式最后一个矩阵的左下角应为 $-1$。按正文所写的 $+1$ 相乘，右端的左下角为 $2\pi^f$；改成 $-1$ 后，乘积才是左边的上三角矩阵。}
\[
\begin{pmatrix}1&\pi^{-f}\\0&1\end{pmatrix}
=\begin{pmatrix}\pi^{-f}&0\\0&\pi^f\end{pmatrix}
\begin{pmatrix}1&0\\\pi^{-f}&1\end{pmatrix}
\begin{pmatrix}\pi^f&1\\1&0\end{pmatrix}.
\]

利用这个恒等式，可以计算差值
\[
\deg t_1^- -\deg t_2^-=
\deg t_1^+-\deg t_2^++2\lambda(\boldsymbol n^+(\boldsymbol y)),
\]
其中 $\lambda(\boldsymbol n^+(\boldsymbol y))=
\sum_{\mathfrak q'\ne\mathfrak p,\mathfrak q}\lambda_{\mathfrak q'}(n_{\mathfrak q'}^+)$，并且
\[
\lambda_{\mathfrak q'}(n_{\mathfrak q'}^+)
=\begin{cases}0,&n_{\mathfrak q'}^+\in\mathcal O_{\mathfrak q'},\\
\deg_{\mathfrak q'}(n_{\mathfrak q'}^+),&n_{\mathfrak q'}^+\notin\mathcal O_{\mathfrak q'}.
\end{cases}
\]
也就是说，若 $n_{\mathfrak q'}^+=\varepsilon_{\mathfrak q'}\pi_{\mathfrak q'}^{-f_{\mathfrak q'}}$，
$f_{\mathfrak q'}>0$，则 $\lambda_{\mathfrak q'}(n_{\mathfrak q'}^+)=-f_{\mathfrak q'}$。于是
\[
\boldsymbol x_0\boldsymbol y=
\boldsymbol x_0\boldsymbol t^+(\boldsymbol y)\boldsymbol n^+(\boldsymbol y)
\boldsymbol k'(\boldsymbol y)
\]
给出 $U(\bar\F_q)$ 中的点，当且仅当
$\deg t_1^+>\deg t_2^+$，且
$2\lambda(\boldsymbol n^+(\boldsymbol y))\geq\deg t_2^+-\deg t_1^+$。\footnote{\textbf{译者注（本次校订）：}这里的次数和截断不等式须与2.1的定义统一。按通常归一化 $|t|=q^{-\deg t}$，应有 $\deg L=\log_q|t_1^+|$；非强分解条件是 $|t_1^+|\leq|t_2^+|$ 且 $|t_2^-|<|t_1^-|$。令 $A=\deg t_1^+-\deg t_2^+$，则结合上面的差值公式，等价条件为 $A\geq0$ 且 $A+2\lambda(n)<0$。此外，全局次数中每个局部赋值须乘以该素位的次数，所以 $\lambda_{\mathfrak q'}(\varepsilon\pi^{-f})=-f\deg(\mathfrak q')$；只假定 $\mathfrak p,\mathfrak q$ 为有理点，并不使其余素位也都具有次数1。}


引入群
\[
T^{(1)}(A)=\left\{\begin{pmatrix}t_1&0\\0&t_2\end{pmatrix}\ \middle|\ |t_1/t_2|=1\right\}.
\]
由简单计数，分解轨道对 $\operatorname{Fr}^m$ 的迹的贡献为
\[
\sum_{a\in T(K)/K}
2\operatorname{vol}(T^{(1)}(A)/(T(K)Z(A)))\cdot\widetilde I'_D(a)\delta(m,a),
\]
其中\footnote{\textbf{译者注（本次校订）：}下式的被积函数中，第一个 $k'$ 应为 $(k')^{-1}$，即取 $\chi'_D((k')^{-1}n^{-1}ank')$。这是对 $n^{-1}an$ 作共轭，也与紧接着的 Selberg 积分中的 $k^{-1}n^{-1}ank$ 一致。}
\[
\widetilde I'_D(a)=-\int_{\mathcal K'\times N^+(A_{\mathfrak p,\mathfrak q})}
\chi'_D(k'n^{-1}ank')\lambda(n)\,dk'\,dn.
\]

% 原书第379页
再将它与 Selberg 迹公式中的对应项比较，后者中的积分为
\[
\int_{\mathcal K\times N^+(A)}
\chi_D\circ\Phi_m^{(\mathfrak p)}\circ\Phi_m^{(\mathfrak q)}
(k^{-1}n^{-1}ank)\,\widetilde\lambda(n)\,dk\,dn.
\]
比较这两个积分时，需要注意第二个积分还包含变量 $n_{\mathfrak p},n_{\mathfrak q}$ 积分，且
\[
\widetilde\lambda(n)=
\sum_{\mathfrak q'\ne\mathfrak p,\mathfrak q}\lambda_{\mathfrak q'}(n_{\mathfrak q'})
+\lambda_{\mathfrak p}(n_{\mathfrak p})+\lambda_{\mathfrak q}(n_{\mathfrak q})
=\lambda(n')+\lambda_{\mathfrak p}(n_{\mathfrak p})+\lambda_{\mathfrak q}(n_{\mathfrak q}).
\]
最后只须证明下面的积分
\[
\int\chi_D\circ\Phi_m^{(\mathfrak p)}\circ\Phi_m^{(\mathfrak q)}
(k^{-1}n^{-1}ank)\lambda_{\mathfrak p}(n_{\mathfrak p})\,dk\,dn
\]
为零，将 $\mathfrak p$ 换为 $\mathfrak q$ 也一样。首先注意，只有
$\operatorname{div}(a)=\pm m\mathfrak p-(\pm m\mathfrak q)$ 时，才可能有非零贡献。这很清楚，因为 $a$ 在 $\mathfrak p,\mathfrak q$ 之外均须为单位。由对
$\Theta_{\mathfrak p}(a,\Phi_m^{(\mathfrak p)})$ 与
$\Theta_{\mathfrak q}(a,\Phi_m^{(\mathfrak q)})$ 的已知结果，就得到这一断言。

若条件成立，则下面的积分消失：
\[
\int_{N^+(K_{\mathfrak p})}
\Phi_m^{(\mathfrak p)}(n_{\mathfrak p}^{-1}an_{\mathfrak p})
\lambda_{\mathfrak p}(n_{\mathfrak p})\,dn_{\mathfrak p}.
\]
这是因为被积函数本身为零。原文由
$\lambda_{\mathfrak p}(n_{\mathfrak p})=0$ 推出
$n_{\mathfrak p}\notin\mathcal O_{\mathfrak p}$，然后易验证
\[
n_{\mathfrak p}^{-1}an_{\mathfrak p}
=an_{\mathfrak p}^{-a}n_{\mathfrak p}
\notin\operatorname{supp}(\Phi_m^{(\mathfrak p)}).
\]
\footnote{译注：所供扫描页此处显示 $\lambda_{\mathfrak p}(n_{\mathfrak p})=0$，但上一页的定义是当 $n_{\mathfrak p}\in\mathcal O_{\mathfrak p}$ 时 $\lambda_{\mathfrak p}=0$。此处的等号疑应为不等号；论证要处理的是权函数非零的情形。}
这也可由 $\Phi_m^{(\mathfrak p)}$ 与 $f_n$ 的关系式直接验证。因此，Selberg 迹公式中该积分的值为
\[
I_D(a)\cdot\Theta_{\mathfrak p}(a,\Phi_m^{(\mathfrak p)})
\cdot\theta_{\mathfrak q}(a,\Phi_m^{(\mathfrak q)}).
\]
Selberg 迹公式中这些项之前有因子 $\frac12$，但
$\Theta_{\mathfrak p}(a,\Phi_m^{(\mathfrak p)})\theta_{\mathfrak q}(a,\Phi_m^{(\mathfrak q)})$
取值为一的频度，是 $\delta(a,m)$ 的四倍。因此两侧的分解贡献也相等，主要等式得证。
\begin{thebibliography}{9}
\bibitem[1]{c39-1} S. Gelbart and H. Jacquet, \textit{Forms on GL(2) from the analytic point of view}, 本论文集，第1册，第213--251页。
\bibitem[2]{c39-2} H. Jacquet and R. Langlands, \textit{Automorphic forms on GL(2)}, Lecture Notes in Math., vol. 114, Springer-Verlag, New York, 1970.
\bibitem[3]{c39-3} R. Langlands, \textit{Modular forms and $\ell$-adic representations. II}, Proc. Summer School Antwerp, Lecture Notes 349, 1973, pp. 361--500.
\bibitem[4]{c39-4} R. Langlands, \textit{Base change for GL(2)}, 油印讲义，Inst. for Advanced Study, Princeton, N.J., 1975.
\bibitem[5]{c39-5} R. Langlands, \textit{On the zeta-functions of some Shimura varieties}, 油印讲义，Inst. for Advanced Study, Princeton, N.J., 1977.
\bibitem[6]{c39-6} D. Mumford, 致 Harder 与 Kazhdan 的未发表信件。
\bibitem[7]{c39-7} H. Saito, \textit{Automorphic forms and algebraic extensions of number fields}, Lecture Notes in Math., Kyoto University, 8, 1975.
\bibitem[8]{c39-8} P. Deligne, \textit{La conjecture de Weil}, Inst. Hautes Études Sci. Publ. Math. 43 (1974), 273--307.
\bibitem[9]{c39-9} P. Deligne, \textit{Poids dans la cohomologie des variétés algébriques}, Actes du Congrès Internat. des Math., Vancouver, 1974.
\end{thebibliography}
\begin{flushleft}\small 波恩数学研究所（Mathematisches Institut, Bonn）\\哈佛大学（Harvard University）\end{flushleft}
